\documentclass[12pt]{amsart}
\usepackage[T1]{fontenc}
\usepackage{lmodern}
\input{glyphtounicode}
\pdfgentounicode=1
\usepackage{amsmath,amssymb,mathtools}
\usepackage[table]{xcolor}
\usepackage{array}
\usepackage{booktabs}
\usepackage[margin=1in]{geometry}
\usepackage{url}
\usepackage{natbib}
\usepackage[pdfauthor={311859884+mt0-svg@users.noreply.github.com},
            pdftitle={The constant of doubly uniform regret in one-dimensional online linear regression is three},
            pdfkeywords={online linear regression, minimax regret, doubly uniform regret, self-normalized martingale},
            bookmarksnumbered=true]{hyperref}
\hypersetup{colorlinks=true, linkcolor=blue, citecolor=blue, urlcolor=blue}
\newcommand{\repo}{https://github.com/mt0-svg/regret-kappa}
% robust, so that amsart does not uppercase the address in the author line
\DeclareRobustCommand{\authorlink}{\href{\repo/issues/new/choose}{\nolinkurl{311859884+mt0-svg@users.noreply.github.com}}}

\renewcommand{\baselinestretch}{1.15}
\setlength{\emergencystretch}{3em}

\newtheorem{theorem}{Theorem}[section]
\newtheorem{lemma}[theorem]{Lemma}
\newtheorem{proposition}[theorem]{Proposition}
\newtheorem{corollary}[theorem]{Corollary}
\theoremstyle{definition}
\newtheorem{definition}[theorem]{Definition}
\newtheorem{remark}[theorem]{Remark}
\newtheorem{computation}[theorem]{Computation}
\newtheorem{question}[theorem]{Question}
\numberwithin{equation}{section}

\newcommand{\R}{\mathbb{R}}
\newcommand{\E}{\mathbb{E}}
\newcommand{\Prob}{\mathbb{P}}
\newcommand{\Reg}{\mathbf{Reg}}
\newcommand{\Regs}{\mathbf{Reg}^{\star}}
\newcommand{\clip}{\operatorname{clip}}
\newcommand{\sgn}{\operatorname{sgn}}
\newcommand{\Ghat}{\widehat{\Gamma}}
\newcommand{\yh}{\hat y}
\newcommand{\Info}{\mathcal{I}}
\newcommand{\Succ}{\mathcal{S}}
\newcommand{\ch}{\operatorname{ch}}
% Lean declarations in typewriter, breakable after an underscore or a dot.
\ExplSyntaxOn
\cs_new:Npn \__lname_char:n #1 { #1 \str_if_in:nnT { _. } {#1} { \allowbreak } }
\NewDocumentCommand \lname { m } { \texttt { \str_map_function:nN {#1} \__lname_char:n } }
\ExplSyntaxOff
% The map from the numbered statements to the formalization and the code. These macros print
% nothing; placed right after a \label, they are read from main.tex to write the table of the
% companion repository's README.
%   \leanmap{A, B}        the Lean declarations of the statement, full names
%   \codemap{S > O, ...}  the scripts S with their recorded outputs O (paths from the root of the repository)
%   \mapnote{text}        how the formal or computed result differs from the statement
\newcommand{\leanmap}[1]{\ignorespaces}
\newcommand{\codemap}[1]{\ignorespaces}
\newcommand{\mapnote}[1]{\ignorespaces}
% A result the formalization or a computation has not delivered yet: shown in the draft, and a
% build with \UseFinal defined stops at the first one left.
\newif\iffinal
\ifdefined\UseFinal\finaltrue\fi
\definecolor{pendingred}{RGB}{176,0,32}
\newcommand{\pending}[2]{\iffinal\errmessage{Pending item left (#1)}\fi{\leavevmode\color{pendingred}[pending, #1: #2]}}

\begin{document}

\title[Doubly uniform regret in dimension one]{The constant of doubly uniform regret in one-dimensional online linear regression is three}
\author{\authorlink}

\begin{abstract}
We consider online linear regression in dimension one when the features are revealed one at a time, with no bound known in advance, and the regret is measured against every linear predictor, with no bound on its parameter. For outcomes bounded by $B$ and a horizon $T\ge355713$, both known to the learner, the minimax regret lies between $B^2(3\log T-2\log\log T-15.2)$ and $B^2(3\log T+0.37)$. The upper bound is the regret of an explicit deterministic learner; the lower bound is forced by an explicit randomized adversary, whose features can be taken in $[0,1]$. Lean~4 proves both bounds. Chen, Qian, Rakhlin and Zhivotovskiy (COLT 2026) showed that this regret is of order $\log T$, which answers a question of Gaillard et al.~(ALT 2019); we find its constant, $3$. With the features known in advance the constant is $1$, so revealing them on the fly triples it.
\end{abstract}

\subjclass[2020]{Primary 68Q32; Secondary 62J05, 62C20, 60G42, 91A05}
\keywords{online linear regression, minimax regret, doubly uniform regret, self-normalized martingale, mixability, van Trees inequality}

\maketitle

\definecolor{leangreen}{RGB}{0,122,61}
\newcommand{\badge}[2]{\colorbox{#1}{\strut\textcolor{white}{\sffamily\bfseries\small #2}}}
\begin{center}
\footnotesize
\renewcommand{\arraystretch}{1.35}
\begin{tabular}{>{\raggedright\arraybackslash}p{0.2\textwidth} >{\raggedright\arraybackslash}p{0.7\textwidth}}
\rowcolor{black!6}
\multicolumn{2}{>{\raggedright\arraybackslash}p{0.93\textwidth}}{\textit{This is AI-generated research: the results, proofs and code were found and written by AI. Credit goes to all the humans whose work it builds on.}}\\
\badge{leangreen}{Proved in Lean 4} & Theorem~\ref{thm:main} with its constants, Theorems~\ref{thm:upper} and~\ref{thm:lower} with the lemmas of Sections~\ref{sec:setting} to~\ref{sec:lower} and Appendices~\ref{app:upper} and~\ref{app:lower}, the numerical steps and the drift certificate of Computation~\ref{comp:drift} included, Corollary~\ref{cor:bounded} and Remark~\ref{rem:rand} are proved in Lean~4 with Mathlib, with no hypothesis and only the standard axioms (\lname{RegretKappa.main}).\\
\rowcolor{black!6}
\multicolumn{2}{>{\raggedright\arraybackslash}p{0.93\textwidth}}{Lean proofs, code, certificates and the \TeX{} source: \href{\repo}{\texttt{github.com/mt0-svg/regret-kappa}}}\\
\end{tabular}
\end{center}


\clearpage
\section{Introduction}\label{sec:intro}

We determine the leading constant of the minimax regret of online linear regression in dimension one when the features are revealed one at a time and the regret is measured against every linear predictor: it is $(3+o(1))B^2\log T$ for outcomes bounded by $B$ and horizon $T$, both known to the learner. The game is the following.
\begin{equation}\label{eq:protocol}
\begin{minipage}{0.84\textwidth}
The learner and the adversary know $B>0$ and the horizon $T\ge1$. For $t=1,\dots,T$:
\begin{enumerate}
\item the adversary reveals a feature $x_t\in\R$;
\item the learner predicts $\yh_t\in\R$;
\item the adversary reveals an outcome $y_t\in[-B,B]$.
\end{enumerate}
\end{minipage}
\end{equation}
No bound is known on $x_t$, and none is put on the comparator $\theta$: on the play $(x,y)=(x_1,\dots,x_T;y_1,\dots,y_T)$ the regret of the learner is
\begin{equation}\label{eq:regret}
\Reg(T)=\sum_{t=1}^T(\yh_t-y_t)^2-\inf_{\theta\in\R}\sum_{t=1}^T(\theta x_t-y_t)^2,
\end{equation}
the regret $\Reg_\Theta(T)$ of \citet[Section~2.1]{CQRZ} with $d=1$ and $\Theta=\R$, and the minimax regret is
\begin{equation}\label{eq:minimax}
\Regs_T(B)=\inf_{\mathcal L}\ \sup_{x\in\R^T,\ y\in[-B,B]^T}\Reg(T),
\end{equation}
the infimum over the deterministic learners $\mathcal L$, which map the past to a prediction (Section~\ref{sec:setting}), taken in the extended real line. A bound on $\Regs_T(B)$ is thus uniform over the comparators and over the features, \emph{doubly uniform} in the terms of \citet{GaillardALT19}.

When the features are known to the learner from the start, the leading constant is $1$: \citet[Theorems~7 and~4]{GaillardALT19} proved that the minimax regret in dimension $d$ is then $dB^2\log T+O(\log\log T)$, with an explicit learner and a lower bound for features in $[0,1]^d$, and they asked whether revealing the features on the fly costs more (Question~\ref{q:GaillardALT19} below). \citet[Theorem~5 and Proposition~6]{CQRZ} showed that $\Regs_T(B)$ is of order $\log T$ (the lower half also follows from \citet[Theorem~4]{GaillardALT19}, since knowing the features in advance can only help the learner), and that in dimension $d\ge2$ the doubly uniform regret is linear in $T$ \citep[Theorem~7]{CQRZ}. Our main result gives the constant, with explicit bounds on both sides.

\begin{theorem}\label{thm:main}
\leanmap{RegretKappa.main, RegretKappa.mainChain, RegretKappa.boundU_le, RegretKappa.kappaEqThree_squeeze}
\mapnote{Lean states the theorem as \lname{RegretKappa.MainTheorem}, with the minimax regret in the extended real numbers; it proves it from \lname{RegretKappa.mainChain}, the form with the bounds $B^2b(T)$ and $B^2u(T)$ of Theorems~\ref{thm:lower} and~\ref{thm:upper} (\lname{RegretKappa.MainTheoremChain}), and \lname{RegretKappa.boundU_le}, $u(T)\le3\log T+0.37$ for $T\ge355713$, and derives the limit from these closed bounds (\lname{RegretKappa.kappaEqThree_squeeze}).}
For every $B>0$ and every integer $T\ge355713$, the minimax regret \eqref{eq:minimax} satisfies
\[
B^2\bigl(3\log T-2\log\log T-15.2\bigr)\le\Regs_T(B)\le B^2\bigl(3\log T+0.37\bigr).
\]
Consequently $\Regs_T(B)/(B^2\log T)\to3$ as $T\to\infty$.
\end{theorem}

The upper bound is the regret $B^2u(T)$ of an explicit learner (Theorem~\ref{thm:upper}), and the lower bound $B^2b(T)$ is forced by an explicit randomized adversary (Theorem~\ref{thm:lower}). The adversary forces the same bound on the expected regret of randomized learners, with features in $[0,1]$ and outcomes in $\{-B,0,B\}$ (Corollary~\ref{cor:bounded} and Remark~\ref{rem:rand}), the features of the lower bound of \citet[Theorem~4]{GaillardALT19} for known features.

\begin{question}[{\citealp[Section~4.1]{GaillardALT19}}]\label{q:GaillardALT19}
``The open question is therefore whether we can improve the lower bound and make it larger for strategies that only discover the features on the fly, or if a doubly uniform regret upper bound of $dB^2\ln(T)+\mathcal O_T(1)$ over the $\mathbf u$ and the $\mathbf x_t,\,y_t$ is also possible in the case of sequentially revealed features.''
\end{question}

Here $\mathbf u\in\R^d$ is the comparator, and the question concerns features with $\|\mathbf x_t\|\le X$ for a bound $X$. \citet[Theorem~5, Proposition~6 and Theorem~7]{CQRZ} answered it, with the orders recalled above. In dimension one, Theorem~\ref{thm:main} adds the constant, and Corollary~\ref{cor:bounded} shows that it is the same for features in $[0,1]$: the leading term is $3B^2\log T$, so no bound of the form $B^2\log T+O(1)$ holds uniformly when the features are revealed on the fly. Our lower bound holds against every learner, whatever it knows; the learner of our upper bound uses $B$ and $T$ (Section~\ref{sec:remarks}). For the lower bound, \citet[Section~4.1]{GaillardALT19} suggest random features that are dependent over time, or features built as functions of the strategy's past predictions. The adversary of Theorem~\ref{thm:lower} does the first: its features depend on the past outcomes, and it never looks at the predictions.

\subsection*{Previous work}
With a scale fixed in advance, the constant is $1$. Against a comparator penalized by a multiple of $\|\theta\|^2$, with features revealed on the fly, the forecaster of \citet{Vovk01} and \citet{AW01} has regret at most $dB^2\log T$ plus a term that depends on the penalty and on a bound on the features \citetext{\citealp[Theorem~1]{Vovk01}; \citealp[Theorem~5.6]{AW01}}, and \citet[Theorem~2]{Vovk01} proved that its leading term $dB^2\log T$ cannot be lowered. \citet[Theorem~15]{MHvE} gave a forecaster that needs no bound on the outcomes, with leading constant $1$ and a bound that grows with $\|\theta\|$ and with the norms of the features. None of these bounds is uniform over the comparators and the features. In dimension $d\ge2$ no bound of order $\log T$ is uniform: besides the linear growth of the square loss regret recalled above, \citet[Theorem~6]{Kot17} showed for linear losses that a regret bound that depends on the comparator only through the norm $(\sum_t(\theta\cdot x_t)^2)^{1/2}$, which is invariant under linear changes of the features, must grow like $\sqrt T$. So dimension one is the case in which a constant exists.

\subsection*{Method}
Both halves work with the state $\rho_t=S_t/\sqrt{V_t}$, where $S_t=\sum_{i\le t}x_iy_i$ and $V_t=\sum_{i\le t}x_i^2$. The regret equals $\sum_t(\yh_t^2-2\yh_ty_t)+\rho_T^2$ (Lemma~\ref{lem:identity}), and each round moves $\rho$ to $\rho\sqrt{1-s}+\varepsilon y\sqrt s$, where $\varepsilon$ is the sign of $x_t$ and $s=x_t^2/V_t$ is a step that the adversary chooses freely (Lemma~\ref{lem:state}). The discrete analysis of \citet[Section~3]{HLPR} for two experts uses the same map, in the variables $\rho\sqrt{1-s}$ and $\sqrt s$, with the step fixed by the round. The proof of \citet[Theorem~4]{CQRZ} runs the map with a step chosen freely by the martingale, and bounds $\E\exp(c\rho_t^2)$ round by round with Hoeffding's lemma.

The learner (Section~\ref{sec:upper}) plays against the potential $2\log(C(\lambda+\Gamma(\rho)))$, where $\lambda$ counts the rounds left and $\Gamma$ mixes the exponentials $e^{\theta\rho}$ over $|\theta|\ge1$. Mixability of the square loss \citetext{\citealp[Section~2, Example~4]{Vovk98}; \citealp[Example~3.13]{HKW}; \citealp[Section~2.4, Lemma~2]{Vovk01}} controls the rounds with small steps, the monotonicity of $\Gamma$ those with large steps, and each round costs an increment $3$ of $\lambda$; at the end the potential must dominate $\rho_T^2$ for $|\rho_T|\le\sqrt T$, which forces $C$ of order $\sqrt T$. The adversary (Section~\ref{sec:lower}) plays in two phases. In the first, of $\lfloor T/2\rfloor$ rounds, it drives $\rho_t^2$ up to a level $L$ close to $2\log T$ with fair random signs as outcomes, against which no prediction gains anything, which leaves the term $\rho_t^2\approx2\log T$ in the regret. In the second, it draws a parameter from a prior centred at the least squares estimate and generates the outcomes from it; by the van Trees inequality \citep[Section~2, inequality~(3)]{GL}, on which the lower bound of \citet[Theorem~4]{GaillardALT19} for known features also rests, estimating this parameter costs every learner about $\log T$ more. The two halves reach $3$ by different splits, as the following heuristic shows; the proofs do not use it. For the learner, $u(T)=2\log\bigl(e^2+C(3T+\Gamma(0))\bigr)$ is about $2\log(3T)+2\log C$, and $2\log C\approx\log T$: the rounds pay $2\log T$ through the increments, and the end pays $\log T$ through $C$. For the adversary, $8e^{\rho^2/2}$ gains at least $1$ per move on average in phase~1 and stays below $8.8\,e^{L/2}$, so the state reaches the level $L$ within about $e^{L/2}$ moves, and $L\approx2\log T$ is the level reachable in $T/2$ rounds; this pays $2\log T$. In phase~2, round $i$ costs the learner about $2/i$, the squared error of estimating the parameter from the outcomes before it, once $i$ exceeds about $\pi\sqrt{T\log T}$; summed over the $T/2$ rounds this is about $\log T-\log\log T$.

\subsection*{Organisation}
Sections~\ref{sec:setting} to~\ref{sec:proofmain} prove Theorem~\ref{thm:main}, with the proofs of the lemmas and the certificates in Appendices~\ref{app:upper} and~\ref{app:lower}. Section~\ref{sec:remarks} gathers open questions, and Section~\ref{sec:checks} describes the formal proof and its data. Logarithms are natural throughout.

\section{The game and the regret identity}\label{sec:setting}

This section fixes the learners and the adversaries of the game \eqref{eq:protocol} and rewrites the regret as a function of one real state, $\rho_t=S_t/\sqrt{V_t}$, which both strategies of the paper read.

\subsection{Learners and adversaries}
Fix $B>0$ and an integer $T\ge1$. A \emph{deterministic learner} $\mathcal L$ is a sequence of maps $\mathcal L_t\colon\R^t\times\R^{t-1}\to\R$, and it predicts $\yh_t=\mathcal L_t(x_1,\dots,x_t;y_1,\dots,y_{t-1})$; its regret on a play is \eqref{eq:regret}, and $\Regs_T(B)$ in \eqref{eq:minimax} is the infimum over these learners. An adaptive adversary chooses $x_t$ and $y_t$ as functions of everything revealed before; against a fixed deterministic learner it produces a single play, so the supremum in \eqref{eq:minimax} is also the supremum over adaptive adversaries. An adversary that ignores the predictions and plays the play $(x^i,y^i)$ with probability $p_i$, for $i$ in a finite set $I$, gives a deterministic learner the expected regret $\sum_ip_i\Reg(x^i,y^i)$. Randomized learners are defined in Section~\ref{sec:proofmain}.

For a play we write, as \citet[Section~2]{CQRZ} do,
\[
S_t=\sum_{i\le t}x_iy_i,\qquad V_t=\sum_{i\le t}x_i^2,\qquad R_t=\frac{S_t^2}{V_t}\quad(R_t=0\text{ if }V_t=0),
\]
and we put $\rho_t=S_t/\sqrt{V_t}$, so that $\rho_t^2=R_t$; further $s_t=x_t^2/V_t\in[0,1]$, and $\varepsilon_t=\sgn x_t$. The conventions for the empty and degenerate cases are $\rho_t=0$ and $s_t=0$ when $V_t=0$, $\varepsilon_t=1$ when $x_t=0$, and $\rho_0=S_0=V_0=0$. Finally $\clip(z)=\max(-1,\min(1,z))$. Since $|\clip(z)-y|\le|z-y|$ for every $y\in[-1,1]$, clipping a prediction to $[-1,1]$ never increases its loss when $|y|\le1$.

\subsection{Scaling and the regret identity}
Both strategies of the paper are described for $B=1$; the next lemma transfers them to every $B$.

\begin{lemma}[scaling]\label{lem:scaling}
\leanmap{RegretKappa.Corollaries.regret_rescale_inv, RegretKappa.Lower.regret_rescale, RegretKappa.Upper.bestLinearLoss_mul}
\mapnote{Lean proves the identity for every $B\ne0$.}
Let $B>0$, let $\mathcal L$ be a deterministic learner, and let $\mathcal L^B$ predict $\yh_t=B\,\mathcal L_t(x_1,\dots,x_t;y_1/B,\dots,y_{t-1}/B)$. For every play $(x,y)$, the regret of $\mathcal L^B$ on $(x,y)$ is $B^2$ times the regret of $\mathcal L$ on $(x,y/B)$.
\end{lemma}

Since $(\mathcal L^{1/B})^B=\mathcal L$, every deterministic learner is of the form $\mathcal L^B$. The learner of Section~\ref{sec:upper} is described for $B=1$ and run for general $B$ in the form $\mathcal L^B$, which uses $B$; the adversary of Section~\ref{sec:lower} is described for $B=1$, and its outcomes are multiplied by $B$ for general $B$. The next identity, the completion of squares of \citet[Section~2.1]{CQRZ} in dimension one, puts the comparator into the state; both halves of the proof start from it.

\begin{lemma}[regret identity]\label{lem:identity}
\leanmap{RegretKappa.Upper.bestLinearLoss_closed, RegretKappa.Upper.sum_sq_sub_bestLinearLoss, RegretKappa.Lower.regret_eq}
For every play, $\inf_{\theta\in\R}\sum_t(\theta x_t-y_t)^2=\sum_ty_t^2-R_T$. Hence
\[
\Reg(T)=\sum_{t=1}^T\bigl(\yh_t^2-2\yh_ty_t\bigr)+\rho_T^2 .
\]
\end{lemma}

\begin{proof}
The sum equals $V_T\theta^2-2S_T\theta+\sum_ty_t^2$. If $V_T>0$, it is minimal at $\theta=S_T/V_T$, with value $\sum_ty_t^2-S_T^2/V_T$. If $V_T=0$, every $x_t$ vanishes, so $S_T=0$ and the sum is the constant $\sum_ty_t^2$. Subtracting from $\sum_t(\yh_t-y_t)^2=\sum_t(\yh_t^2-2\yh_ty_t)+\sum_ty_t^2$ gives the formula, since $\rho_T^2=R_T$. This is \lname{RegretKappa.Lower.regret_eq} in the formalization.
\end{proof}

The next lemma, which the learner and the adversary both use, shows how one round moves the state and what it costs.

\begin{lemma}[the state]\label{lem:state}
\leanmap{RegretKappa.Upper.rho_sq_le, RegretKappa.Upper.sq_sum_mul_le, RegretKappa.UpperSharp.step_sharp}
\mapnote{In Lean the update (b) is the rewriting inside \lname{RegretKappa.UpperSharp.step_sharp}, which also handles $V_{t-1}=0$ and $x_t=0$.}
For every play and every $t\ge1$:
\begin{enumerate}
\item[(a)] if $|y_i|\le1$ for all $i\le t$, then $|\rho_t|\le\sqrt t$;
\item[(b)] $\rho_t=\rho_{t-1}\sqrt{1-s_t}+\varepsilon_ty_t\sqrt{s_t}$;
\item[(c)] if $\yh_t=\varepsilon_t\eta$ with $\eta\in\R$, then $\yh_t^2-2\yh_ty_t=\eta^2-2\eta\,\varepsilon_ty_t$.
\end{enumerate}
\end{lemma}

The proofs of Lemmas~\ref{lem:scaling} and~\ref{lem:state}, in Appendix~\ref{app:upper}, are direct: the substitution $\theta\mapsto\theta/B$ for the first, and the Cauchy-Schwarz inequality and the expansion of $S_t$ for the second. With outcomes in $[-1,1]$, the game is therefore played on $\rho$, as follows.
\begin{equation}\label{eq:reduced}
\begin{minipage}{0.84\textwidth}
In round $t$, from the state $\rho_{t-1}$:
\begin{enumerate}
\item the adversary chooses a step $s=s_t\in[0,1]$;
\item the learner, who knows $s$, chooses $\eta\in\R$ and predicts $\yh_t=\varepsilon_t\eta$;
\item the adversary chooses $\tilde y=\varepsilon_ty_t\in[-1,1]$;
\item the learner pays $\eta^2-2\eta\tilde y$, and the state moves to $\rho_t=\rho_{t-1}\sqrt{1-s}+\tilde y\sqrt s$.
\end{enumerate}
At the end the learner pays $\rho_T^2$.
\end{minipage}
\end{equation}
When $V_{t-1}>0$ the adversary reaches every $s\in[0,1)$, with $x_t^2=V_{t-1}s/(1-s)$, and never $s=1$; when $V_{t-1}=0$ the state is $\rho_{t-1}=0$ and $s_t\in\{0,1\}$.

\section{An explicit learner}\label{sec:upper}

This section proves the upper bound of Theorem~\ref{thm:main} with an explicit learner whose regret is at most $u(T)$ on every play with outcomes in $[-1,1]$. The proofs of its lemmas are in Appendix~\ref{app:upper}.

By the regret identity (Lemma~\ref{lem:identity}), the learner must keep $\sum_t(\yh_t^2-2\yh_ty_t)+\rho_T^2$ small in the game \eqref{eq:reduced}. For the comparator $\theta/\sqrt{V_t}$, half of the gain $\sum_{i\le t}y_i^2-\sum_{i\le t}(\theta x_i/\sqrt{V_t}-y_i)^2$ over the comparator $0$ is $\theta\rho_t-\theta^2/2$, and the learner plays against a mixture of the exponentials $e^{\theta\rho-\theta^2/2}$ of these half gains over $|\theta|\ge1$, as the method of mixtures does for self-normalized sums \citep[Theorem~2.1]{dlPKL}, where the weight is a Gaussian of fixed scale; Section~\ref{subsec:small} explains the cutoff and the weight.

For $|\theta|\ge1$ let $w(\theta)=e^{-\theta^2/2}\bigl(|\theta|^{-1}+2|\theta|^{-3}\bigr)$, and for $\rho\in\R$ let
\begin{equation}\label{eq:Gamma}
\Gamma(\rho)=\int_{|\theta|\ge1}e^{\theta\rho}w(\theta)\,d\theta=2\int_1^\infty\cosh(\theta\rho)\,e^{-\theta^2/2}\Bigl(\frac1\theta+\frac2{\theta^3}\Bigr)d\theta .
\end{equation}
Recall from Section~\ref{sec:setting} that $\rho_t=S_t/\sqrt{V_t}$, $s_t=x_t^2/V_t$ and $\varepsilon_t=\sgn x_t$, where $S_t=\sum_{i\le t}x_iy_i$ and $V_t=\sum_{i\le t}x_i^2$, and that $\clip(z)=\max(-1,\min(1,z))$. For $\lambda>0$, $\rho\in\R$ and $s\in[0,1]$, with $c=\sqrt{1-s}$ and $\sigma=\sqrt s$ here and below, let
\begin{equation}\label{eq:learner}
\eta(\lambda,\rho,s)=\clip\Bigl(\frac12\log\frac{\lambda+\Gamma(\rho c+\sigma)}{\lambda+\Gamma(\rho c-\sigma)}\Bigr),\qquad\yh_t=\varepsilon_t\,\eta\bigl(\lambda_0+3(T-t),\rho_{t-1},s_t\bigr).
\end{equation}
The learner that predicts $\yh_t$ in round $t$ is the \emph{potential learner} with parameters $T$ and $\lambda_0>0$. By Lemma~\ref{lem:state}(b), $\rho c\pm\sigma$ are the states to which the outcomes $\varepsilon_ty_t=\pm1$ lead from $\rho=\rho_{t-1}$, so the prediction is the clipped half log-ratio of the potential at these two states. The increment $3$ of $\lambda$ per round is the larger of the one-round costs $2.8857$ of small steps and $2.7053$ of large steps (Lemmas~\ref{lem:small} and~\ref{lem:large}), rounded up.

\begin{theorem}\label{thm:upper}
\leanmap{RegretKappa.UpperSharp.theoremU, RegretKappa.UpperSharp.theoremULog,
  RegretKappa.Corollaries.regret_rescale_learnerU_le, RegretKappa.boundU_le,
  RegretKappa.UpperSharp.ulog_bound, RegretKappa.UpperSharp.const_c_le}
\mapnote{Lean proves the first form for every $T\ge0$ and the second for $T\ge1$, for outcomes in $[-1,1]$ with the learner written out (\lname{RegretKappa.UpperSharp.learnerU}), the case of $[-B,B]$ with the learner of Lemma~\ref{lem:scaling}, and the bound $u(T)\le3\log T+0.37$ as \lname{RegretKappa.boundU_le}. The bounds $4e^{-1/2}/3\le0.81$, $\sqrt T\ge596$ and $2\log(3/\sqrt{2\pi})\le0.366$ of the proof are steps of \lname{RegretKappa.UpperSharp.ulog_bound} and \lname{RegretKappa.boundU_le}, and $2\sqrt{2\pi}\,e^2/3\le12.35$ is \lname{RegretKappa.UpperSharp.const_c_le}.}
Let $T\ge1$ be an integer, $C=(\sqrt T+1)/\sqrt{2\pi}$ and $\lambda_0=e^2/C$, and let $\Gamma$ be as in \eqref{eq:Gamma}. On every play with outcomes in $[-1,1]$, the potential learner \eqref{eq:learner} with parameters $T$ and $\lambda_0$ has regret \eqref{eq:regret}
\begin{align*}
\Reg(T)\le u(T)&:=2\log\Bigl(e^2+\frac{(\sqrt T+1)(3T+2e^{-1/2})}{\sqrt{2\pi}}\Bigr)\\
&\le3\log T+2\log\frac{3}{\sqrt{2\pi}}+\frac{2}{\sqrt T}+\frac{0.81}{T}+\frac{12.35}{T^{3/2}}.
\end{align*}
In particular $u(T)\le3\log T+0.37$ for $T\ge355713$. For outcomes in $[-B,B]$, the learner obtained from \eqref{eq:learner} by Lemma~\ref{lem:scaling} has $\Reg(T)\le B^2u(T)$ on every play.
\end{theorem}

In $u(T)$, the term $e^2=C\lambda_0$ pays the terminal condition for $|\rho|\le2$, where the bound of Lemma~\ref{lem:terminal} on $\Gamma$ does not apply, which is why $\lambda_0=e^2/C$; the factor $C$ pays it up to $|\rho|=\sqrt T$; $3T$ is the increment $3$ of each of the $T$ rounds; and $2e^{-1/2}=\Gamma(0)$ is the potential at the start $\rho_0=0$. The learner reads the feature $x_t$ only through its sign $\varepsilon_t$ and the step $s_t$, predicts in $[-1,1]$, and each prediction costs two evaluations of $\Gamma$, a one-dimensional integral. By Theorem~\ref{thm:lower}, the leading constant $3$ of $u(T)$ cannot be lowered.

\subsection{Proof of Theorem~\ref{thm:upper}}\label{subsec:proofU}
For $0\le k\le T$ let $\Phi_k(\rho)=2\log\bigl(C(\lambda_0+3k+\Gamma(\rho))\bigr)$, the potential with $k$ rounds left, with $C$ and $\lambda_0$ as in Theorem~\ref{thm:upper}. The proof rests on three statements: the inequality for one round (Proposition~\ref{prop:P}), its sum over the rounds (Lemma~\ref{lem:potential}), and the terminal condition (Lemma~\ref{lem:terminal}). In the first, $\eta(\lambda,\rho,s)$, $c=\sqrt{1-s}$ and $\sigma=\sqrt s$ are as in \eqref{eq:learner}.

\begin{proposition}[one round]\label{prop:P}
\leanmap{RegretKappa.UpperSharp.step_abstract_sharp, RegretKappa.UpperSharp.step_paper}
\mapnote{Lean proves it in the normalization $\Gamma/2$ of the formalization, where the increment is $3/2$, and in the variables of the learner with increment $3$.}
For every $\lambda>0$, $\rho\in\R$, $s\in[0,1]$ and $y\in[-1,1]$, with $\eta=\eta(\lambda,\rho,s)$,
\begin{equation}\label{eq:P}
\eta^2-2\eta y+2\log\bigl(\lambda+\Gamma(\rho c+y\sigma)\bigr)\le2\log\bigl(\lambda+3+\Gamma(\rho)\bigr).
\end{equation}
\end{proposition}

The proof, in Section~\ref{subsec:proofP}, reduces the outcomes to $\pm1$ by convexity, then bounds the cost of a round by $2.8857$ for steps $s\le2/3$, by mixability, and by $2.7053$ for $s\ge2/3$, by the monotonicity of $\Gamma$; the split at $2/3$ keeps both bounds below $3$. Summed over the rounds, the proposition bounds the potential along the play.

\begin{lemma}[the potential]\label{lem:potential}
\leanmap{RegretKappa.UpperSharp.potential_leU}
On every play with outcomes in $[-1,1]$, the potential learner \eqref{eq:learner} with parameters $T$ and $\lambda_0$ satisfies, for $0\le n\le T$,
\[
\sum_{t\le n}\bigl(\yh_t^2-2\yh_ty_t\bigr)+\Phi_{T-n}(\rho_n)\le\Phi_T(0).
\]
\end{lemma}

The last input is the terminal condition: at the end, the potential must dominate the term $\rho_T^2$ of the regret identity.

\begin{lemma}[terminal condition]\label{lem:terminal}
\leanmap{RegretKappa.UpperSharp.terminal_ineq, RegretKappa.UpperSharp.terminal_sharp}
For $\rho\ge2$, $\sqrt{2\pi}\,e^{\rho^2/2}\le(\rho+1)\Gamma(\rho)$, with $\Gamma$ as in \eqref{eq:Gamma}. Consequently, with $C=(\sqrt T+1)/\sqrt{2\pi}$ and $\lambda_0=e^2/C$ as in Theorem~\ref{thm:upper}, $C(\lambda_0+\Gamma(\rho))\ge e^{\rho^2/2}$ for $|\rho|\le\sqrt T$, that is $\Phi_0(\rho)\ge\rho^2$ there.
\end{lemma}

The proof, in Section~\ref{subsec:proofP}, bounds $e^{-\rho^2/2}\Gamma(\rho)$ from below by the integral of $\theta^{-1}e^{-(\theta-\rho)^2/2}$ over $\theta\ge1$, and this integral by the Gaussian tail.

\begin{proof}[Proof of Theorem~\ref{thm:upper}]
By Lemma~\ref{lem:potential} with $n=T$, $\sum_{t\le T}(\yh_t^2-2\yh_ty_t)+\Phi_0(\rho_T)\le\Phi_T(0)$. By Lemma~\ref{lem:state}(a), $|\rho_T|\le\sqrt T$, so $\Phi_0(\rho_T)\ge\rho_T^2$ by Lemma~\ref{lem:terminal}, and Lemma~\ref{lem:identity} gives $\Reg(T)\le\Phi_T(0)=2\log(C\lambda_0+C(3T+\Gamma(0)))=u(T)$, since $C\lambda_0=e^2$ and $\Gamma(0)=2e^{-1/2}$ (Lemma~\ref{lem:Gamma}(a)). For the second form write $C(3T+\Gamma(0))=\frac{3T^{3/2}}{\sqrt{2\pi}}(1+T^{-1/2})(1+\Gamma(0)/(3T))$, where $2\log\frac{3T^{3/2}}{\sqrt{2\pi}}=3\log T+2\log\frac3{\sqrt{2\pi}}$, and use $\log(1+v)\le v$ three times:
\[
u(T)-3\log T-2\log\frac3{\sqrt{2\pi}}\le\frac2{\sqrt T}+\frac{2\Gamma(0)}{3T}+\frac{2\sqrt{2\pi}\,e^2}{3T^{3/2}}=\frac2{\sqrt T}+\frac{4e^{-1/2}}{3}\,\frac1T+\frac{2\sqrt{2\pi}\,e^2}{3}\,\frac1{T^{3/2}},
\]
with $4e^{-1/2}/3\le0.81$ and $2\sqrt{2\pi}\,e^2/3\le12.35$ (Computation~\ref{comp:upper}). For $T\ge355713$, $\sqrt T\ge596$, so that $2/\sqrt T<0.00336$ and $0.81/T+12.35/T^{3/2}<10^{-5}$; and $2\log(3/\sqrt{2\pi})=\log(9/(2\pi))\le0.366$, since $9/(2\pi)\le1+0.366+0.366^2/2\le e^{0.366}$. Hence $u(T)\le3\log T+0.37$. For outcomes in $[-B,B]$, by Lemma~\ref{lem:scaling} the regret of the learner $\mathcal L^B$ on a play $(x,y)$ is $B^2$ times the regret of \eqref{eq:learner} on $(x,y/B)$, whose outcomes lie in $[-1,1]$. The bound $\Reg(T)\le u(T)$ is \lname{RegretKappa.UpperSharp.theoremU} in the formalization.
\end{proof}

\section{An explicit adversary}\label{sec:lower}

This section proves the lower bound of Theorem~\ref{thm:main} with the two-phase adversary outlined in the introduction. The proofs of its lemmas and the drift certificate are in Appendix~\ref{app:lower}. Until part (c) of Theorem~\ref{thm:lower}, $B=1$, and the outcomes of the adversary lie in $\{-1,0,1\}$.


\subsection{The statement}\label{subsec:statement}
For an integer $T\ge3$ let
\begin{gather}
j_0=\lceil\log(3\log T)\rceil,\qquad L=2\log\frac{T}{20e\,j_0},\qquad k_0=T-\lfloor T/2\rfloor,\notag\\
c_1=\frac12+\frac43-\frac8{\pi^2}+\log(\pi^2)+\frac{\pi^2/4-2}{\pi^2e^2}\le3.33,\notag\\
b(T)=\bigl(1-e^{-j_0}\bigr)\Bigl[L+\log\frac{k_0}{(\sqrt L+2)^2}-c_1-\frac{\log k_0}{2k_0}-\frac1{k_0}\Bigr].\label{eq:bT}
\end{gather}
The level $L$ is chosen so that $j_0$ blocks of about $8.8e\cdot e^{L/2}$ rounds take $0.44\,T$ of the $\lfloor T/2\rfloor$ rounds of phase~1 (condition (C2) of Section~\ref{subsec:prooflower}), and $j_0$ so that phase~1 fails with probability at most $e^{-j_0}\le1/(3\log T)$, which costs at most $1$ in $b(T)$. In the bracket, $L$ is the term $\rho^2$ of the regret identity at the end of phase~1, $\log(k_0/(\sqrt L+2)^2)$ the cost of estimating the parameter of phase~2 from the state $|\rho|=\sqrt L$, and $c_1$ collects the constants of phase~2, among them the second moment $\frac43-\frac8{\pi^2}$ and the Fisher information $\pi^2/4$ of the prior (Lemma~\ref{lem:prior}).


\begin{theorem}\label{thm:lower}
\leanmap{RegretKappa.LowerSharp.lowerSharpAdv, RegretKappa.LowerSharp.lowerSharpSimplified,
  RegretKappa.LowerSharp.lowerSharpBound, RegretKappa.LowerSharp.c1_le, RegretKappa.LowerSharp.B0_ge,
  RegretKappa.LowerSharp.bT_ge, RegretKappa.LowerSharp.simplified_two, RegretKappa.LowerSharp.B0_nonneg}
\mapnote{(a), (b) and (c) are \lname{RegretKappa.LowerSharp.LowerSharpAdv}, \lname{RegretKappa.LowerSharp.LowerSharpSimplified} and \lname{RegretKappa.LowerSharp.LowerSharpBound}; the proof of (a) builds the plays of $\mathcal A_T$ (\lname{RegretKappa.LowerSharp.adv_core}). The bound $c_1\le3.33$ is \lname{RegretKappa.LowerSharp.c1_le}; the proof of (b) is \lname{RegretKappa.LowerSharp.B0_ge}, \lname{RegretKappa.LowerSharp.bT_ge}, \lname{RegretKappa.LowerSharp.simplified_two} and \lname{RegretKappa.LowerSharp.B0_nonneg}, with the steps \lname{RegretKappa.LowerSharp.two_log_a_add_two_le}, \lname{RegretKappa.LowerSharp.small_k0} and \lname{RegretKappa.LowerSharp.const_le}.}
Let $T\ge T_0=355713$ be an integer, and let $b(T)$ be as in \eqref{eq:bT}.
\begin{enumerate}
\item[(a)] There is an adversary that ignores the predictions and plays finitely many plays $(x^i,y^i)$, with probabilities $p_i$, with features $x^i_t\ge0$ and outcomes $y^i_t\in\{-1,0,1\}$, against which every deterministic learner has $\sum_ip_i\Reg(x^i,y^i)\ge b(T)$.
\item[(b)] $b(T)\ge3\log T-\log\log T-2\log(\log\log T+2.1)-14.6\ge3\log T-2\log\log T-15.2$.
\item[(c)] For every $B>0$, every deterministic learner has a play with outcomes in $[-B,B]$ on which $\Reg(T)\ge B^2b(T)$.
\end{enumerate}
\end{theorem}

The proof shows that the adversary $\mathcal A_T$ of Section~\ref{subsec:adversary} is one as in (a). The threshold $T_0$ is set so that $L\ge14.5$ for every $T\ge T_0$, the value that condition (C1) of Section~\ref{subsec:prooflower} asks of $L$ (Lemma~\ref{lem:side}). The leading term of $b(T)$ is $3\log T$, that of Theorem~\ref{thm:upper}.

\subsection{The adversary}\label{subsec:adversary}
Let $T_1=\lfloor T/2\rfloor$, the length of phase~1, so that $k_0=T-T_1$ is the length of phase~2. The adversary $\mathcal A_T$ uses independent random signs $\xi_1,\dots,\xi_{T_1}$, uniform on $\{-1,1\}$, a real random variable $Z$, independent of the signs, with density $g^2$ on $[-2,2]$, where
\[
g(z)=\frac{\cos(\pi z/4)}{\sqrt2},
\]
and further independent coin flips for the outcomes of phase~2. The variable $Z$ is the error of the prior of phase~2 in the scale of the state. Its support is bounded, so that the outcome probabilities of phase~2 lie in $[0,1]$, and its density vanishes at $\pm2$ and has finite Fisher information, as the van Trees inequality (Lemma~\ref{lem:vT}) requires.


\emph{The move rule.} In phase~1 the adversary chooses the step as a function $s(\rho)$ of the state. A step proportional to $e^{-\rho^2/2}$ makes $8e^{\rho^2/2}$ gain a constant in expectation at each move: for a small step $s$, a move with a fair sign raises $e^{\rho^2/2}$ on average by about $s\,e^{\rho^2/2}/2$ (Lemma~\ref{lem:drift}). Let $s\colon\R\to(0,1)$ be the even function equal to $s_j$ on $\{|\rho|\in[j/10,(j+1)/10)\}$ for $j=0,\dots,27$, and to $\frac58e^{-\rho^2/2}$ for $|\rho|\ge2.8$, where $10^4(s_0,\dots,s_{27})$ is
\[
\begin{gathered}
(6242,6180,6058,5879,5648,5373,5060,4718,4355,3980,3601,3226,2861,2513,\\
2184,1880,1602,1352,1129,934,764,620,497,395,311,242,187,142),
\end{gathered}
\]
the values of $\frac58e^{-m^2/2}$ at the midpoints $m=j/10+1/20$, rounded to four decimals. The factor $5/8$, the cells of width $1/10$ and the four decimals are one admissible choice, the one that Computation~\ref{comp:drift} certifies. The stopped move is $s_L(\rho)=s(\rho)$ if $\rho^2<L$, and $s_L(\rho)=0$ if $\rho^2\ge L$.

\emph{Phase 1, rounds $1$ to $T_1$.} Round $1$ is $x_1=1$, $y_1=\xi_1$. For $1\le t<T_1$, round $t+1$ is
\[
x_{t+1}=\sqrt{V_t\,s_L(\rho_t)/(1-s_L(\rho_t))},\qquad y_{t+1}=\xi_{t+1}.
\]
Once $\rho_t^2\ge L$ the features are $0$, so $\rho$ and $V$ keep their values. Let $\Succ=\{\rho_{T_1}^2\ge L\}$ be the event that phase~1 succeeds.

\emph{Phase 2, rounds $T_1+1$ to $T$, on $\Succ$.} Let $r=|\rho_{T_1}|+2$ and $\theta^*=(\rho_{T_1}+Z)/\sqrt{V_{T_1}}$, that is $S_{T_1}/V_{T_1}+Z/\sqrt{V_{T_1}}$: a comparator drawn from a prior centred at the least squares estimate of phase~1. For $i=1,\dots,k_0$, round $T_1+i$ is
\[
x_{T_1+i}=\alpha_i\sqrt{V_{T_1}},\quad\alpha_i=\frac1{\sqrt2\,r}\sqrt{\frac i{k_0}},\qquad y_{T_1+i}=\begin{cases}1&\text{with probability }(1+\theta^*x_{T_1+i})/2,\\-1&\text{otherwise,}\end{cases}
\]
the outcomes being independent given $\theta^*$ and the past. Since $\theta^*x_{T_1+i}=(\rho_{T_1}+Z)\alpha_i$ and $|Z|\le2$, we have $|\theta^*x_{T_1+i}|\le\nu_i:=\sqrt{i/(2k_0)}\le1/\sqrt2$, so these are probabilities, and $y_{T_1+i}$ has conditional mean $\theta^*x_{T_1+i}$. The squared features $\alpha_i^2$ grow linearly in $i$, which keeps the factors $1-\nu_i^2$ of Lemma~\ref{lem:phase2}, from the variance of the outcomes, from costing more than a constant (a heuristic, which the proofs do not use).
Phase~2 follows the scheme of the proof of \citet[Theorem~4]{GaillardALT19}, after \citet[Theorem~2]{Vovk01}: a prior on the parameter, outcomes drawn from it, and the van Trees inequality for each round; here the prior is centred at the least squares estimate of phase~1.


\emph{Failure, off $\Succ$}: $x_t=y_t=0$ for $t>T_1$.

The adversary never looks at the predictions. Its phase-1 features are functions of the earlier signs and its phase-2 features of $(\rho_{T_1},V_{T_1})$, so they carry no information on $Z$; and since the signs and the phase-2 outcomes take finitely many values, $\mathcal A_T$ plays finitely many plays, each with its probability. The next lemma identifies phase~1 with a Markov chain on the state, for the proofs of Lemmas~\ref{lem:phase1} and~\ref{lem:split}.

\begin{lemma}[the chain]\label{lem:chain}
\leanmap{RegretKappa.LowerSharp.sum_xS_sq, RegretKappa.LowerSharp.sum_xS_mul, RegretKappa.LowerSharp.rhoS_eq_chainOf}
For $1\le t<T_1$, with $V_t=\sum_{i\le t}x_i^2$ and $S_t=\sum_{i\le t}x_iy_i$ as in Section~\ref{sec:setting}, $V_{t+1}=V_t/(1-s_L(\rho_t))$ and $\rho_{t+1}=\rho_t\sqrt{1-s_L(\rho_t)}+\xi_{t+1}\sqrt{s_L(\rho_t)}$, and $S_t=\rho_t\sqrt{V_t}$ with $V_t\ge1$ for $1\le t\le T_1$. So $(\rho_t)_{1\le t\le T_1}$ is the Markov chain on $\R$ started at $\rho_1=\xi_1$ that moves from $\rho$ to $\rho\sqrt{1-s_L(\rho)}+\sqrt{s_L(\rho)}$ or to $\rho\sqrt{1-s_L(\rho)}-\sqrt{s_L(\rho)}$, with probability $1/2$ each; we call it the stopped chain, since it stays at $\rho$ once $\rho^2\ge L$.
\end{lemma}

\subsection{Proof of Theorem~\ref{thm:lower}}\label{subsec:prooflower}
Put $a=\sqrt L$ and $a_+=a+\sqrt{5/8}\,e^{-L/4}$, the largest value of $|\rho|$ that phase~1 can reach (Lemma~\ref{lem:overshoot}). On $\Succ$ write $\rho=\rho_{T_1}$, $Q=\sum_{i\le k_0}\alpha_i^2$, and $\yh'_i=\yh_{T_1+i}$, $y'_i=y_{T_1+i}$ for $1\le i\le k_0$. The proof uses three side conditions on $T$,
\begin{enumerate}
\item[(C1)] $L\ge14.5$;
\item[(C2)] $j_0(8.8\,e\,e^{L/2}+1)\le T_1-1$;
\item[(C3)] $k_0\ge\pi^2e^2(a_++2)^2$,
\end{enumerate}
and four lemmas: the side conditions hold (Lemma~\ref{lem:side}), phase~1 succeeds with high probability and the learner gains nothing in it (Lemma~\ref{lem:phase1}), the regret splits on each play (Lemma~\ref{lem:split}), and phase~2 costs the learner at least a function $\beta$ of the state (Lemma~\ref{lem:phase2}). Their proofs are in Appendix~\ref{app:lower}, with that of Lemma~\ref{lem:chain}, which follows from the definition of the features and Lemma~\ref{lem:state}(b).

\begin{lemma}\label{lem:side}
\leanmap{RegretKappa.LowerSharp.level_ge, RegretKappa.LowerSharp.C2, RegretKappa.LowerSharp.C3,
  RegretKappa.LowerSharp.exp_33_4_le, RegretKappa.LowerSharp.five_case, RegretKappa.LowerSharp.j0_le}
Conditions (C1), (C2) and (C3) hold for every integer $T\ge T_0=355713$.
\end{lemma}

Condition (C1) follows from the definition of $j_0$, (C2) from the choice of $L$, since $8.8\,e\,e^{L/2}j_0=0.44\,T$, and (C3) from $(a_++2)^2\le4\log T+18$.

Phase~1 drives $\rho^2$ to the level $L$ with high probability, and the learner gains nothing in it, since its outcomes are fair signs, as in the lower bounds of \citet[Section~2.1]{CQRZ}.

\begin{lemma}[phase 1]\label{lem:phase1}
\leanmap{RegretKappa.LowerSharp.prob_fail_le, RegretKappa.LowerSharp.abs_rhoS_le, RegretKappa.LowerSharp.avg_P1_nonneg}
Recall that $j_0=\lceil\log(3\log T)\rceil$, $T_1=\lfloor T/2\rfloor$ and $\Succ=\{\rho_{T_1}^2\ge L\}$. Assume (C1) and (C2). Then:
\begin{enumerate}
\item[(1)] $\Prob(\Succ)\ge1-e^{-j_0}$;
\item[(2)] $|\rho_{T_1}|\le a_+$, and $|\rho_{T_1}|\ge a$ on $\Succ$;
\item[(3)] for every deterministic learner, $\E\sum_{t\le T_1}(\yh_t^2-2\yh_ty_t)\ge0$.
\end{enumerate}
\end{lemma}

Part (1) rests on a drift function: the move rule makes $U(\rho)=8e^{\rho^2/2}$ gain at least $1$ in expectation at each move, while $U$ stays below $8.8\,e^{L/2}$ on the states the chain reaches, which gives a geometric tail for the time the chain takes to reach the level $L$ (Section~\ref{subsec:phase1}).

On a play in $\Succ$, the regret splits into the cost of phase~1, the term $\rho^2$ of the regret identity, and a term $\Delta(z)$ for phase~2. For $z=Z$, $\Delta(Z)$ is the regret over phase~2 against the comparator $\theta^*$, minus the penalty $V_{T_1}(\theta^*-S_{T_1}/V_{T_1})^2=Z^2$, plus the gap between $\theta^*$ and the final least squares estimate: the past acts as a prior.

\begin{lemma}[the regret on a play]\label{lem:split}
\leanmap{RegretKappa.LowerSharp.regret_succ, RegretKappa.LowerSharp.regret_fail, RegretKappa.LowerSharp.comparator_eq}
On a play of $\mathcal A_T$ in $\Succ$, with $\rho=\rho_{T_1}$, $\alpha_i=\sqrt{i/k_0}/(\sqrt2\,r)$ and $r=|\rho|+2$ as in Section~\ref{subsec:adversary}, for every $z\in\R$,
\[
\Reg(T)=\sum_{t\le T_1}(\yh_t^2-2\yh_ty_t)+\rho^2+\Delta(z),
\]
\[
\Delta(z)=-z^2+\sum_{i=1}^{k_0}\Bigl[(\yh'_i-y'_i)^2-\bigl((\rho+z)\alpha_i-y'_i\bigr)^2\Bigr]+(1+Q)\Bigl(\frac{\rho+\sum_i\alpha_iy'_i}{1+Q}-\rho-z\Bigr)^2 .
\]
On a play off $\Succ$, $\Reg(T)\ge\sum_{t\le T_1}(\yh_t^2-2\yh_ty_t)$.
\end{lemma}

The proof expands the squares, with $S_{T_1}=\rho\sqrt{V_{T_1}}$, and applies Lemma~\ref{lem:identity}.

The van Trees inequality bounds the expectation of $\Delta(Z)$ from below through the information that the outcomes of phase~2 carry on $Z$.

\begin{lemma}[phase 2]\label{lem:phase2}
\leanmap{RegretKappa.LowerSharp.phase2S, RegretKappa.LowerSharp.beta_ge}
\mapnote{Lean proves the first claim for every radius $r\ge|\rho|+2$; in part (iv) of the proof it bounds $\sum_{1\le i<k_0}\log(i/k_0)$ through $k_0^{k_0}/k_0!\le e^{k_0}$ in place of the integral of $\log$.}
Let $T_1=\lfloor T/2\rfloor$ and $k_0=T-T_1$ be the lengths of phases~1 and~2, and $Z$, of density $g^2$, the variable of Section~\ref{subsec:adversary}. Fix phase-1 signs $\xi_1,\dots,\xi_{T_1}$ on $\Succ$, and let $\rho=\rho_{T_1}$, $r=|\rho|+2$, $\alpha_i=\sqrt{i/k_0}/(\sqrt2\,r)$, $\nu_i=\sqrt{i/(2k_0)}$ and $\Delta$ be as above, $J=\pi^2/4$, and $c_1$ as in \eqref{eq:bT}. Let $\Info_0=J$ and $\Info_i=\Info_{i-1}+\alpha_i^2/(1-\nu_i^2)$ for $1\le i\le k_0$. For every deterministic learner, $\E[\Delta(Z)\mid\xi_1,\dots,\xi_{T_1}]\ge\beta(\rho)$, where
\[
\beta(\rho)=\sum_{i=1}^{k_0}\frac{\alpha_i^2}{\Info_{i-1}}-\E Z^2+\frac{1+Q}{\Info_{k_0}};
\]
and if $k_0\ge\pi^2e^2r^2$, then $\beta(\rho)\ge\underline\beta(\rho):=\log(k_0/r^2)-c_1-\log(k_0)/(2k_0)-1/k_0$.
\end{lemma}

The first claim applies the van Trees inequality (Lemma~\ref{lem:vT}) to each prediction of phase~2 and to the final least squares estimate. The second bounds $\sum_i\alpha_i^2/\Info_{i-1}$ from below by a summation by parts and an integral of $\log$, and $(1+Q)/\Info_{k_0}$ from below with $4Q\ge\pi^2e^2$. The term $\log(k_0/r^2)$ of $\underline\beta$, about $\log T$ since $k_0\approx T/2$ and $r^2\approx2\log T$, is what phase~2 adds to the term $\rho_{T_1}^2\ge L\approx2\log T$ of phase~1.

\begin{proof}[Proof of Theorem~\ref{thm:lower}]
(a) Fix a deterministic learner. By Lemma~\ref{lem:split} with $z=Z$, $\Reg(T)=\sum_{t\le T_1}(\yh_t^2-2\yh_ty_t)+\rho_{T_1}^2+\Delta(Z)$ on $\Succ$, and $\Reg(T)\ge\sum_{t\le T_1}(\yh_t^2-2\yh_ty_t)$ off $\Succ$. Conditioning on the phase-1 signs and applying Lemma~\ref{lem:phase2} on $\Succ$,
\[
\E\,\Reg(T)\ge\E\sum_{t\le T_1}(\yh_t^2-2\yh_ty_t)+\E\bigl[\mathbf 1_\Succ\bigl(\rho_{T_1}^2+\beta(\rho_{T_1})\bigr)\bigr],
\]
and the first expectation is nonnegative by Lemma~\ref{lem:phase1}(3), which applies by (C1) and (C2) (Lemma~\ref{lem:side}). On $\Succ$, $a\le|\rho_{T_1}|\le a_+$ by Lemma~\ref{lem:phase1}(2), so $r\le a_++2$ and (C3) makes the bound $\underline\beta$ of Lemma~\ref{lem:phase2} valid. The function $v\mapsto v^2+\log(k_0/(v+2)^2)$ is nondecreasing on $[1,\infty)$, with derivative $2v-2/(v+2)>0$, and $a\ge1$; hence on $\Succ$,
\[
\rho_{T_1}^2+\underline\beta(\rho_{T_1})\ge b_0:=L+\log\frac{k_0}{(a+2)^2}-c_1-\frac{\log k_0}{2k_0}-\frac1{k_0},
\]
and, since $b_0\ge0$ (Appendix~\ref{app:lower}), $\E\,\Reg(T)\ge\Prob(\Succ)\,b_0\ge(1-e^{-j_0})\,b_0=b(T)$ by Lemma~\ref{lem:phase1}(1). The expectation is the weighted sum over the finitely many plays of $\mathcal A_T$, whose features are nonnegative and whose outcomes lie in $\{-1,0,1\}$, which proves (a).

(b) Collecting the terms of $b_0$, with $j_0\le\log\log T+2.1$ and $L\le2\log T$, gives $b_0\ge3\log T-\log\log T-2\log(\log\log T+2.1)-13.6$, and $e^{-j_0}\le1/(3\log T)$ costs at most $1$ more; the proof is in Appendix~\ref{app:lower}.

(c) Let $B>0$ and let $\mathcal L$ be a deterministic learner. By (a), applied to the learner $\mathcal L^{1/B}$ of Lemma~\ref{lem:scaling}, some play $(x^i,y^i)$ has $\Reg(x^i,y^i)\ge b(T)$ for $\mathcal L^{1/B}$, since a weighted average with weights of sum $1$ is at most the largest term. Since $(\mathcal L^{1/B})^B=\mathcal L$, Lemma~\ref{lem:scaling} gives that the regret of $\mathcal L$ on $(x^i,By^i)$ is $B^2$ times that, at least $B^2b(T)$, and the outcomes $By^i_t$ lie in $[-B,B]$. This is \lname{RegretKappa.LowerSharp.lowerSharpBound} in the formalization.
\end{proof}

\section{Proof of Theorem~\ref{thm:main}; bounded features and randomized learners}\label{sec:proofmain}

\begin{proof}[Proof of Theorem~\ref{thm:main}]
Let $B>0$ and $T\ge355713$. By Theorem~\ref{thm:lower}(c), every deterministic learner has a play with outcomes in $[-B,B]$ and $\Reg(T)\ge B^2b(T)$, so the supremum in \eqref{eq:minimax} is at least $B^2b(T)$ for every learner, and $\Regs_T(B)\ge B^2b(T)\ge B^2(3\log T-2\log\log T-15.2)$ by Theorem~\ref{thm:lower}(b). The learner of Theorem~\ref{thm:upper} for outcomes in $[-B,B]$ has $\Reg(T)\le B^2u(T)$ on every play, and $u(T)\le3\log T+0.37$ by the same theorem, so $\Regs_T(B)\le B^2(3\log T+0.37)$. Finally $(3\log T-2\log\log T-15.2)/\log T\to3$ and $(3\log T+0.37)/\log T\to3$; dividing by $B^2\log T$ gives the limit. This is \lname{RegretKappa.main} in the formalization.
\end{proof}

\subsection{Bounded features and randomized learners}
The adversary $\mathcal A_T$ of Section~\ref{subsec:adversary} plays features that grow geometrically in phase~1. Since its plays are finitely many, its features can be scaled into $[0,1]$ at no cost, and since it ignores the predictions, its bound extends to randomized learners, which we now define.

A \emph{randomized learner} is a probability space $(\Omega,\mathcal F,\mu)$ with a deterministic learner $\mathcal L^\omega$ for each $\omega\in\Omega$, such that for each round $t$ and each history $(x_1,\dots,x_t;y_1,\dots,y_{t-1})$ the prediction $\mathcal L^\omega_t(x_1,\dots,x_t;y_1,\dots,y_{t-1})$ is a measurable function of $\omega$. Let $\Reg^\omega(T)$ be the regret of $\mathcal L^\omega$ on a play. Since the comparator $\theta=0$ has loss $\sum_ty_t^2$, $\Reg^\omega(T)+\sum_ty_t^2\ge0$, and the expected regret on the play is
\[
\E\,\Reg(T)=\int_\Omega\Bigl(\Reg^\omega(T)+\sum_ty_t^2\Bigr)\mu(d\omega)-\sum_ty_t^2\in(-\infty,+\infty],
\]
the integral of a nonnegative measurable function, with no integrability assumed. A deterministic learner is a randomized learner whose family is constant, and its expected regret is its regret. An adversary that ignores the predictions and plays the play $(x^i,y^i)$ with probability $p_i$, for $i$ in a finite set $I$, is \emph{oblivious}: its plays do not depend on the learner's randomness. Against it the expected regret is defined in the same way, by integrating $\sum_ip_i(\Reg^\omega(x^i,y^i)+\sum_t(y^i_t)^2)$ against $\mu$ and subtracting $\sum_ip_i\sum_t(y^i_t)^2$; against a deterministic learner it is $\sum_ip_i\Reg(x^i,y^i)$, as in Section~\ref{sec:setting}. For a set $\mathcal Y\subseteq[-B,B]$ let $\Regs_T(B;[0,1],\mathcal Y)$ be the minimax regret \eqref{eq:minimax} with the supremum restricted to the plays with features in $[0,1]$ and outcomes in $\mathcal Y$.

\begin{corollary}\label{cor:bounded}
\leanmap{RegretKappa.mainBoundedFeatures, RegretKappa.Corollaries.boundedFeatures_of_mixture, RegretKappa.Corollaries.le_advRegret}
\mapnote{Lean states it as \lname{RegretKappa.MainBoundedFeatures}, with the two minimax regrets \lname{RegretKappa.Corollaries.minimaxRegretBF} and \lname{RegretKappa.Corollaries.minimaxRegretBFI} in the extended real numbers, for randomized learners on probability spaces in every universe. The sharper form of the sentence after the corollary, with $B^2b(T)$ and $B^2u(T)$, is \lname{RegretKappa.mainBoundedFeaturesChain} (\lname{RegretKappa.MainBoundedFeaturesChain}).}
Let $B>0$ and let $T\ge355713$ be an integer. There is an oblivious adversary that plays finitely many plays $(x^i,y^i)$, with probabilities $p_i$, with features in $[0,1]$ and outcomes in $\{-B,0,B\}$, against which every randomized learner has expected regret at least $B^2(3\log T-2\log\log T-15.2)$. For deterministic learners,
\begin{align*}
B^2\bigl(3\log T-2\log\log T-15.2\bigr)&\le\Regs_T\bigl(B;[0,1],\{-B,0,B\}\bigr)\\
&\le\Regs_T\bigl(B;[0,1],[-B,B]\bigr)\le B^2\bigl(3\log T+0.37\bigr).
\end{align*}
\end{corollary}

The proof gives the sharper bounds of Theorems~\ref{thm:lower} and~\ref{thm:upper}: the adversary forces $B^2b(T)$, and $B^2b(T)\le\Regs_T(B;[0,1],\{-B,0,B\})\le\Regs_T(B;[0,1],[-B,B])\le B^2u(T)$, with $b(T)$ as in \eqref{eq:bT} and $u(T)$ as in Theorem~\ref{thm:upper}. Its proof is in Appendix~\ref{app:lower}.

\begin{remark}[randomized learners on one play]\label{rem:rand}
\leanmap{RegretKappa.randLowerBound, RegretKappa.CorollariesCheck.V42b.not_randLowerBound}
\mapnote{Lean states it as \lname{RegretKappa.Corollaries.RandLowerBound} and proves it from Corollary~\ref{cor:bounded}; the reading with the Bochner integral and no measurability is refuted by \lname{RegretKappa.CorollariesCheck.V42b.not_randLowerBound}, with the learner of \lname{RegretKappa.CorollariesCheck.integral_famBot}.}
For every $B>0$ and $\gamma>0$ there is $T_\gamma$ such that for every $T\ge T_\gamma$, every randomized learner has a play with outcomes in $[-B,B]$ on which its expected regret is at least $(3-\gamma)B^2\log T$. By the measurability in the definition of a randomized learner, its expected regret against an oblivious adversary is the weighted average of its expected regrets on the plays, so this follows from Corollary~\ref{cor:bounded}, as $3\log T-2\log\log T-15.2\ge(3-\gamma)\log T$ for large $T$. The measurability cannot be dropped: without it the expected regret is not defined, and with the Bochner integral of the regret in its place, which is $0$ for a function that is not integrable, the statement is false, since two deterministic learners on a two-point space with the trivial $\sigma$-algebra have expected regret $0$ on every play with outcomes in $[-1,1]$. The first statement is \lname{RegretKappa.randLowerBound} in the formalization.
\end{remark}

When the features are known to the learner from the start, the minimax regret in dimension one lies between $B^2(\log T-\log\log T-3)$ for $T\ge8$, with features in $[0,1]$ \citep[Theorem~4]{GaillardALT19}, and $B^2(\log(1+T)+1)$, for all real features \citep[Theorem~7]{GaillardALT19}. By Corollary~\ref{cor:bounded}, revealing the features one at a time raises the constant from $1$ to $3$.

\section{Further remarks}\label{sec:remarks}

\emph{The additive constants.} Theorem~\ref{thm:main} leaves a gap of order $\log\log T$ between the two bounds. In the lower bound, the term $-\log\log T$ of Theorem~\ref{thm:lower}(b) comes from the radius $r=|\rho_{T_1}|+2$ of phase~2, through $\log(k_0/r^2)$ in Lemma~\ref{lem:phase2} with $r^2\approx2\log T$, and the term $-2\log(\log\log T+2.1)$ from the factor $j_0$ in the level $L$. A second order term of order $\log\log T$ is not specific to features revealed on the fly: for the constant feature $x_t=1$, known in advance, the minimax regret is $B^2(\log T-\log\log T+O(\log\log T/\log T))$ \citep[Corollary~1 and Lemma~3, for their loss, which is half of ours]{TW00}. Whether $\Regs_T(1)-3\log T$ is bounded is open. Our guess is that the lower bound is the one to move: a phase~1 that reached the level $L$ in a time of order $e^{L/2}/\sqrt L$, as a diffusion does, in place of $e^{L/2}$, would gain about $\log L$ in the level and cancel the term $-\log\log T$ of the radius. What blocks this route is that a drift function of order $e^{\rho^2/2}/|\rho|$, for a move rule with smaller steps at high levels, has a relative drift that becomes very small as $|\rho|$ grows, too small at high levels for bounds on intervals as in Computation~\ref{comp:drift} to certify.

\emph{Unknown $B$ or $T$.} The learner of Theorem~\ref{thm:upper} uses $T$ through the count $\lambda$ of the rounds left and the scale $C\approx\sqrt{T/(2\pi)}$ of the terminal condition, and $B$ through the rescaling of Lemma~\ref{lem:scaling}. Whether some learner that knows neither $B$ nor $T$ has, for every $T$ and every play of length $T$ with $|y_t|\le B$, regret at most $(3+\varepsilon_T)B^2\log T$, with $\varepsilon_T\to0$ independent of the play and of $B$, is open, and we have no guess. The learner of Theorem~\ref{thm:upper} has this regret but knows $B$ and $T$. The algorithm of \citet[Theorem~5]{CQRZ} has regret of order $B^2\log T$ on every play, with a constant that it does not make explicit, and it also takes the horizon $T$ as input. The lower bound of Theorem~\ref{thm:lower} holds for every learner, so it applies unchanged.

\section{The formal proof and its data}\label{sec:checks}
This section describes how the proof is checked: what the formal proof proves and what it trusts, and which program wrote the data it checks.

Lean~4.34.1 \citep{Lean4} with Mathlib \citep{Mathlib} at its release v4.34.1 proves Theorem~\ref{thm:main}, Theorems~\ref{thm:upper} and~\ref{thm:lower}, Corollary~\ref{cor:bounded} and Remark~\ref{rem:rand}, with the constants of the paper. Their formal statements rest on definitions that transcribe Sections~\ref{sec:intro}, \ref{sec:setting} and~\ref{sec:proofmain}: a deterministic learner for the horizon $T$ is a sequence of maps from $x_1,\dots,x_t$ and $y_1,\dots,y_{t-1}$ to a prediction, the regret is \eqref{eq:regret}, the minimax regret \eqref{eq:minimax} is an infimum of suprema in the extended real numbers, and randomized learners and their expected regret are those of Section~\ref{sec:proofmain}, with the expectation in the extended real numbers. The theorems have no hypothesis and depend on the three standard axioms of Lean and Mathlib, propositional extensionality, the axiom of choice and the soundness of quotients, and on no other axiom. No proof has a function evaluated by compiled code, so no compiled code belongs to the trusted base; the Lean kernel checks every declaration. The numerical steps are inequalities between explicit numbers, or inequalities in one real or integer variable proved by monotonicity or by a polynomial bound, all derived from rational bounds proved in Lean on the constants involved ($e$, $\pi$, square roots, logarithms and $K_0$); the kernel evaluates the rational numbers of the drift certificate of Computation~\ref{comp:drift}, one cell per declaration. Comparator \citep{Comparator} checks ten of these theorems, those of Theorem~\ref{thm:main} and Corollary~\ref{cor:bounded}, each in its closed form and with the bounds of Theorems~\ref{thm:upper} and~\ref{thm:lower}, and those of Theorems~\ref{thm:upper} and~\ref{thm:lower} and of Remark~\ref{rem:rand}: it checks that they prove the statements of separate files, which a program writes from the statement files by copying their definitions and stating the ten theorems without proof, and which import only Mathlib and one another, with no axiom other than these three, and it has every declaration they depend on, those of Mathlib included, checked again by nanoda \citep{nanoda}, an implementation of the Lean kernel written separately. The formal proof follows the paper, with three differences of form: the upper bound works with $\Gamma/2$ in place of $\Gamma$, so with the increment $3/2$; Lemma~\ref{lem:small} is proved from polynomial bounds on $\cosh$ with a remainder, in place of the series of $\cosh$, and these bounds lead to the same quartic $f$; and in part (iv) of the proof of Lemma~\ref{lem:phase2}, $\sum_{1\le i<k_0}\log(i/k_0)$ is bounded through $k_0^{k_0}/k_0!\le e^{k_0}$ in place of the integral of $\log$.

The subdivisions of Computation~\ref{comp:drift} were written into the Lean module by a PARI/GP program \citep{PARI}. The proof does not trust this program, since the kernel checks what it writes.

The computations used PARI/GP~2.17.4 and Lean~4.34.1 with Mathlib~v4.34.1, on one machine, an AMD Ryzen 9 5900X (12 cores, 24 threads), 64~GB of memory. Building the formal proof takes a few minutes on four cores once Mathlib is built, and the check with Comparator about $2$ minutes. The repository \citep{repo} maps every numbered statement to its Lean declarations.

\emph{Data availability.} The repository \citep{repo} holds the source of this paper, the Lean package with the records of its checks, and the program that wrote the subdivisions of the drift certificate.

\section*{Contact}
Questions, corrections and comments are welcome as issues on the repository: \url{\repo/issues/new/choose}.

\appendix

\section{Proofs of Sections~\ref{sec:setting} and~\ref{sec:upper}}\label{app:upper}

This appendix proves Lemmas~\ref{lem:scaling}, \ref{lem:state}, \ref{lem:potential} and~\ref{lem:terminal} and Proposition~\ref{prop:P}, from the lemmas on $\Gamma$ below. As in Section~\ref{sec:upper}, $\Gamma$ is the potential \eqref{eq:Gamma}, $c=\sqrt{1-s}$ and $\sigma=\sqrt s$.

\subsection{Proofs of Lemmas~\ref{lem:scaling}, \ref{lem:state} and~\ref{lem:potential}}

\begin{proof}[Proof of Lemma~\ref{lem:scaling}]
Put $y'=y/B$ and let $\yh'_t$ be the predictions of $\mathcal L$ on $(x,y')$, so that $\mathcal L^B$ predicts $B\yh'_t$. Then $(B\yh'_t-y_t)^2=B^2(\yh'_t-y'_t)^2$ and $\sum_t(\theta x_t-y_t)^2=B^2\sum_t((\theta/B)x_t-y'_t)^2$; since $\theta\mapsto\theta/B$ is a bijection of $\R$, the infimum over $\theta$ is multiplied by $B^2$ as well. This is \lname{RegretKappa.Corollaries.regret_rescale_inv} in the formalization.
\end{proof}

\begin{proof}[Proof of Lemma~\ref{lem:state}]
(a) By the Cauchy-Schwarz inequality, $S_t^2\le(\sum_{i\le t}y_i^2)V_t\le tV_t$. (b) If $V_t=0$, then $x_t=0$ and $V_{t-1}=0$, so $\rho_t=\rho_{t-1}=0$ and $s_t=0$. If $V_t>0$, then $\rho_t=(S_{t-1}+x_ty_t)/\sqrt{V_t}$; the first part is $\rho_{t-1}\sqrt{V_{t-1}/V_t}=\rho_{t-1}\sqrt{1-s_t}$ (both sides vanish when $V_{t-1}=0$), and the second is $\varepsilon_ty_t\sqrt{s_t}$. (c) Use $\varepsilon_t^2=1$. Part (a) is \lname{RegretKappa.Upper.rho_sq_le} in the formalization, where (b) and (c) are steps of the proof of Proposition~\ref{prop:P}.
\end{proof}

\begin{proof}[Proof of Lemma~\ref{lem:potential}]
By induction on $n$. For $n=0$ both sides are $\Phi_T(0)$, since $\rho_0=0$. Let $1\le t\le T$, $\lambda=\lambda_0+3(T-t)>0$, $s=s_t$ and $y=\varepsilon_ty_t\in[-1,1]$. By Lemma~\ref{lem:state}(b) and (c),
\[
\yh_t^2-2\yh_ty_t+\Phi_{T-t}(\rho_t)=\eta^2-2\eta y+2\log C+2\log\bigl(\lambda+\Gamma(\rho_{t-1}c+y\sigma)\bigr)
\]
with $\eta=\eta(\lambda,\rho_{t-1},s)$, and Proposition~\ref{prop:P} bounds the right side by $2\log C+2\log(\lambda+3+\Gamma(\rho_{t-1}))=\Phi_{T-t+1}(\rho_{t-1})$. Adding $\sum_{i<t}(\yh_i^2-2\yh_iy_i)$ and the induction hypothesis at $n=t-1$ gives the claim at $n=t$. This is \lname{RegretKappa.UpperSharp.potential_leU} in the formalization.
\end{proof}

\subsection{The potential}\label{subsec:Gamma}
The proofs of this appendix use the following properties of $\Gamma$: its symmetry and monotonicity (a), its moments (b), for the series of Lemma~\ref{lem:small}, and a bound on its increments (c), (d), for the large steps.

\begin{lemma}\label{lem:Gamma}
\leanmap{RegretKappa.UpperSharp.Gam_neg,
  RegretKappa.UpperSharp.Gam_nonneg,
  RegretKappa.Upper.G_mono,
  RegretKappa.UpperSharp.Gam_zero,
  RegretKappa.UpperSharp.mom_odd,
  RegretKappa.UpperSharp.Mw_one,
  RegretKappa.UpperSharp.Mw_add_two,
  RegretKappa.UpperSharp.Mw_eq_GamC,
  RegretKappa.UpperSharp.Dg_hasSum,
  RegretKappa.UpperSharp.Dg_nonneg,
  RegretKappa.UpperSharp.Dg_mono,
  RegretKappa.UpperSharp.coeff_le,
  RegretKappa.UpperSharp.G_sqrt_sub_le}
\mapnote{Lean proves (c) and (d) in the normalization $\Gamma/2$, with $\delta/2$; it writes $K_n$ as $e^{-1/2}$ times an integer given by the recursion of (b).}
Let $\Gamma$ be as in \eqref{eq:Gamma}.
\begin{enumerate}
\item[(a)] $\Gamma$ is even and nonnegative, $\Gamma(\rho)\le\Gamma(\rho')$ whenever $|\rho|\le|\rho'|$, and $\Gamma(0)=2e^{-1/2}$.
\item[(b)] For $n\ge0$ let $\Gamma_n=2\int_1^\infty\theta^{2n}e^{-\theta^2/2}(\theta^{-1}+2\theta^{-3})\,d\theta$ and $K_n=\int_1^\infty\theta^{2n-1}e^{-\theta^2/2}\,d\theta$. Then $K_1=e^{-1/2}$, $K_{n+1}=e^{-1/2}+2nK_n$ for $n\ge1$, and $\Gamma_n=2(K_n+2K_{n-1})$ for $n\ge1$.
\item[(c)] For $v>0$ let $\delta(v)=\int_1^\infty\theta\sinh(\theta\sqrt v)\,v^{-1/2}\,e^{-\theta^2/2}(\theta^{-1}+2\theta^{-3})\,d\theta$. Then $\delta(v)=\sum_{n\ge1}n\Gamma_nv^{n-1}/(2n)!$, the function $\delta$ is nonnegative and nondecreasing on $(0,\infty)$, and $n\Gamma_n/(2n)!\le2/n!$ for $n\ge1$.
\item[(d)] For $0\le v\le v'$ with $v'>0$, $\Gamma(\sqrt{v'})-\Gamma(\sqrt v)\le(v'-v)\,\delta(v')$.
\end{enumerate}
\end{lemma}

\begin{proof}
(a) The weight $w$ is even, positive and has Gaussian decay, and $\cosh(\theta\rho)$ is even in $\rho$ and nondecreasing in $|\rho|$, so $\Gamma$, written with $\cosh$ as in \eqref{eq:Gamma}, is finite, even, nonnegative and nondecreasing in $|\rho|$. Since $\frac{d}{d\theta}\bigl(-\theta^{-2}e^{-\theta^2/2}\bigr)=e^{-\theta^2/2}(\theta^{-1}+2\theta^{-3})$, $\Gamma(0)=2e^{-1/2}$.

(b) Expanding the bracket gives $\Gamma_n=2(K_n+2K_{n-1})$ for $n\ge1$. Next $K_1=\int_1^\infty\theta e^{-\theta^2/2}d\theta=e^{-1/2}$, and an integration by parts of $\theta^{2n}\cdot\theta e^{-\theta^2/2}$ gives $K_{n+1}=e^{-1/2}+2n\int_1^\infty\theta^{2n-1}e^{-\theta^2/2}d\theta=e^{-1/2}+2nK_n$.

(c) For $\theta\ge1$, $\theta\sinh(\theta\sqrt v)/\sqrt v=\sum_{n\ge1}2n\,\theta^{2n}v^{n-1}/(2n)!$, with nonnegative terms; integrating term by term against $e^{-\theta^2/2}(\theta^{-1}+2\theta^{-3})$ gives the series, whose coefficients are nonnegative, so $\delta$ is nonnegative and nondecreasing. For the bound, the change of variable from $\theta$ to $\theta^2/2$ gives $K_n\le\int_0^\infty\theta^{2n-1}e^{-\theta^2/2}d\theta=2^{n-1}(n-1)!$ for $n\ge1$. For $n\ge2$ this gives $\Gamma_n\le2\bigl(2^{n-1}(n-1)!+2^{n-1}(n-2)!\bigr)\le2^{n+1}(n-1)!$, and $\Gamma_1=2e^{-1/2}+4K_0\le4$, since $K_0\le0.28$ (Computation~\ref{comp:upper}). Since $(2n)!=\binom{2n}{n}(n!)^2\ge2^n(n!)^2$, we get $n\Gamma_n/(2n)!\le n2^{n+1}(n-1)!/(2^n(n!)^2)=2/n!$.

(d) For $\theta\ge1$ the function $v\mapsto\cosh(\theta\sqrt v)=\sum_n\theta^{2n}v^n/(2n)!$ is convex on $[0,\infty)$, with derivative $\theta\sinh(\theta\sqrt v)/(2\sqrt v)$ at $v>0$, so
\[
\cosh(\theta\sqrt{v'})-\cosh(\theta\sqrt v)\le(v'-v)\,\frac{\theta\sinh(\theta\sqrt{v'})}{2\sqrt{v'}}.
\]
Multiplying by $2e^{-\theta^2/2}(\theta^{-1}+2\theta^{-3})$ and integrating over $\theta\ge1$ gives the claim. Part (d), the bound used for the large steps, is \lname{RegretKappa.UpperSharp.G_sqrt_sub_le} in the formalization.
\end{proof}

\subsection{Two outcomes}\label{subsec:two}
The first step of the proof of Proposition~\ref{prop:P} reduces the outcomes to $\pm1$. For $y\in[-1,1]$ write $A(y)=\lambda+\Gamma(\rho c+y\sigma)$, with $c=\sqrt{1-s}$ and $\sigma=\sqrt s$, and $A_\pm=A(\pm1)$. The left side of \eqref{eq:P} is $2\log\bigl(e^{\eta^2/2-\eta y}A(y)\bigr)$.

\begin{lemma}[endpoints]\label{lem:endpoints}
\leanmap{RegretKappa.Upper.G_convex, RegretKappa.Upper.exp_mul_le}
\mapnote{Lean proves the inequality for $\Gamma/2$ in place of $\Gamma$ and for every $\rho c$ and $\sigma$ in $\R$, and for the term $\lambda e^{-\eta y}$ separately.}
For $\eta\in\R$, $\lambda\ge0$, $\rho\in\R$, $s\in[0,1]$ and $y\in[-1,1]$,
\[
e^{-\eta y}A(y)\le\frac{1+y}2\,e^{-\eta}A_++\frac{1-y}2\,e^{\eta}A_- .
\]
In particular the left side of \eqref{eq:P} is at most the larger of its values at $y=1$ and $y=-1$.
\end{lemma}

\begin{proof}
By \eqref{eq:Gamma}, $e^{-\eta y}A(y)=\lambda e^{-\eta y}+\int_{|\theta|\ge1}e^{\theta\rho c}\,e^{y(\theta\sigma-\eta)}\,w(\theta)\,d\theta$. For fixed $\theta$ the functions $y\mapsto e^{-\eta y}$ and $y\mapsto e^{y(\theta\sigma-\eta)}$ are convex, so on $[-1,1]$ each lies below its chord; integrating against the positive weight gives the inequality. Its right side is a convex combination of the values of $e^{-\eta y}A(y)$ at $y=\pm1$. This is \lname{RegretKappa.Upper.G_convex} in the formalization.
\end{proof}

\subsection{Mixability}\label{subsec:mix}
The prediction rests on Hoeffding's lemma for the law on $\{-1,1\}$ with mean $\eta$, in the following form.

\begin{lemma}[two points]\label{lem:mix}
\leanmap{RegretKappa.Upper.two_point, RegretKappa.Upper.two_point_mgf}
\mapnote{Lean derives it from Hoeffding's lemma in Mathlib, applied to the law on two points of \lname{RegretKappa.Upper.two_point_mgf}.}
For $\eta\in[-1,1]$ and $z\in\R$,
\[
(1+\eta)\,e^{z-\eta}+(1-\eta)\,e^{-(z-\eta)}\le2\,e^{(z^2-\eta^2)/2}.
\]
\end{lemma}

\begin{proof}
Let $\kappa(v)=\log\bigl(((1+\eta)e^v+(1-\eta)e^{-v})/2\bigr)$, the logarithm of the moment generating function of the law on $\{-1,1\}$ with mean $\eta$. Then $\kappa(0)=0$, $\kappa'(0)=\eta$, and $\kappa''=1-(\kappa')^2\le1$, since $\kappa'(v)=((1+\eta)e^v-(1-\eta)e^{-v})/((1+\eta)e^v+(1-\eta)e^{-v})$. By Taylor's formula $\kappa(v)\le\eta v+v^2/2$, which is Hoeffding's lemma for this law. At $v=z-\eta$, $\eta v+v^2/2=(z^2-\eta^2)/2$. This is \lname{RegretKappa.Upper.two_point} in the formalization.
\end{proof}

The next lemma turns an inequality of this form into a bound on the cost of the prediction $\eta(\lambda,\rho,s)$ of \eqref{eq:learner} against both outcomes.

\begin{lemma}[the prediction]\label{lem:choice}
\leanmap{RegretKappa.Upper.mix_alg}
\mapnote{Lean proves it for all positive numbers in place of $A_+$ and $A_-$.}
Let $\lambda>0$, $\rho\in\R$ and $s\in[0,1]$, let $A_\pm$ be as in Section~\ref{subsec:two}, and let $W\in\R$ satisfy $(1+\eta)e^{-\eta}A_++(1-\eta)e^{\eta}A_-\le2e^{-\eta^2/2}W$ for every $\eta\in[-1,1]$. Then $\eta^*=\eta(\lambda,\rho,s)$ satisfies
\[
e^{-\eta^*}A_+\le e^{-\eta^{*2}/2}W\qquad\text{and}\qquad e^{\eta^*}A_-\le e^{-\eta^{*2}/2}W .
\]
\end{lemma}

\begin{proof}
The numbers $A_\pm$ are positive, since $\lambda>0$ and $\Gamma\ge0$ (Lemma~\ref{lem:Gamma}(a)). Let $\eta_0=\frac12\log(A_+/A_-)$, so that $\eta^*=\clip(\eta_0)$. If $|\eta_0|\le1$, then $\eta^*=\eta_0$ and $e^{-\eta_0}A_+=e^{\eta_0}A_-$, so the hypothesis at $\eta=\eta_0$ reads $2e^{\eta_0}A_-\le2e^{-\eta_0^2/2}W$, which is both inequalities. If $\eta_0>1$, then $\eta^*=1$, the hypothesis at $\eta=1$ reads $2e^{-1}A_+\le2e^{-1/2}W$, and $e\,A_-<e^{-1}A_+$ since $A_+>e^2A_-$. The case $\eta_0<-1$ is symmetric. This is \lname{RegretKappa.Upper.mix_alg} in the formalization.
\end{proof}

Integrated against a positive measure, Lemma~\ref{lem:mix} gives the hypothesis of Lemma~\ref{lem:choice} for a mixture of exponentials. The two lemmas together are a form of the mixability of the square loss at rate $1/2$ for the outcomes $\pm1$ \citep[Section~2.4, Lemma~2]{Vovk01}; for outcomes in $\{0,1\}$ it is the mixability at rate $2$ of \citet[Section~2, Example~4]{Vovk98} and \citet[Example~3.13]{HKW}, after the change of outcome $y\mapsto2y-1$. The next corollary takes two choices of $W$, one for small steps and one for large steps.

\begin{corollary}\label{cor:mix}
\leanmap{RegretKappa.Upper.G_tangent, RegretKappa.Upper.mix_small, RegretKappa.Upper.mix_large}
\mapnote{Lean proves (a) and (b) in the normalization $\Gamma/2$.}
Let $\lambda>0$, $\rho\in\R$, $s\in[0,1]$, $y\in[-1,1]$, and $\eta=\eta(\lambda,\rho,s)$, with $\eta$, $c=\sqrt{1-s}$ and $\sigma=\sqrt s$ as in \eqref{eq:learner} and $w(\theta)=e^{-\theta^2/2}(|\theta|^{-1}+2|\theta|^{-3})$ as in \eqref{eq:Gamma}.
\begin{enumerate}
\item[(a)] If $s<1$, the left side of \eqref{eq:P} is at most $2\log(\lambda+\Ghat(\rho,s))$, where
\[
\Ghat(\rho,s)=\int_{|\theta|\ge1}e^{\theta^2s/2}\,e^{\theta c\rho}\,w(\theta)\,d\theta<\infty .
\]
\item[(b)] The left side of \eqref{eq:P} is at most $2\log\bigl(\lambda+\max(\Gamma(\rho c+\sigma),\Gamma(\rho c-\sigma))\bigr)$.
\end{enumerate}
\end{corollary}

\begin{proof}
By Lemma~\ref{lem:endpoints} it suffices to take $y=\pm1$, and by Lemma~\ref{lem:choice}, the left side of \eqref{eq:P} at $y=\pm1$ is at most $2\log W$ for every $W$ that satisfies its hypothesis. (a) Since $s<1$, $e^{\theta^2s/2}w(\theta)$ has Gaussian decay and $\Ghat(\rho,s)$ is finite. Multiply Lemma~\ref{lem:mix} at $z=\theta\sigma$ by $e^{\theta c\rho}w(\theta)$, integrate over $|\theta|\ge1$, and add $\lambda$ times Lemma~\ref{lem:mix} at $z=0$; this is the hypothesis with $W=\lambda+\Ghat(\rho,s)$. (b) Lemma~\ref{lem:mix} at $z=0$ gives $(1+\eta)e^{-\eta}+(1-\eta)e^{\eta}\le2e^{-\eta^2/2}$, and $1\pm\eta\ge0$; multiplied by $W=\max(A_+,A_-)$ this gives the hypothesis. Part (a) is \lname{RegretKappa.Upper.mix_small} in the formalization.
\end{proof}

\subsection{Small steps}\label{subsec:small}
For small steps, Corollary~\ref{cor:mix}(a) leaves the comparison of $\Ghat$ with $\Gamma$. The change of variable from $\theta$ to $c\theta$ brings $\Ghat$ to the form of $\Gamma$, with the cutoff moved from $1$ to $c$ and the term $2\theta^{-3}$ multiplied by $c^2$. This is what the weight $w$ of \eqref{eq:Gamma} is made for. The weight $e^{-\theta^2/2}|\theta|^{-1}$ on all of $\R$ would make $\Ghat$ equal to the mixture before the step, since the change of variable turns one into the other (Lemma~\ref{lem:excess}), but it is not integrable at $\theta=0$. Cutting it off at $|\theta|=1$ leaves a positive excess, of order $s\,e^{|\rho|}$, from the moved cutoff; with this weight alone the excess is unbounded in $\rho$, since it is at least $2e^{-1/2}\cosh(\rho\sqrt{1-s})\log(1/\sqrt{1-s})$ for every $s\in(0,1)$. The factor $1+2\theta^{-2}$, which decreases in $|\theta|$, adds the negative term $-4sH(\rho)$ below, which grows like $s\,e^{\rho^2/2}|\rho|^{-3}$ and beats the excess. Its coefficient $2$ makes the weight the derivative of $-\theta^{-2}e^{-\theta^2/2}$, which gives $\Gamma(0)=2e^{-1/2}$ and the moments of Lemma~\ref{lem:Gamma}(b).

\begin{lemma}[the substitution]\label{lem:excess}
\leanmap{RegretKappa.Upper.Ghat_subst}
\mapnote{Lean proves it in the normalization $\Gamma/2$.}
For $s\in[0,1)$ and $\rho\in\R$, with $c=\sqrt{1-s}$ and $\Ghat$ as in Corollary~\ref{cor:mix},
\[
\Ghat(\rho,s)=2\int_c^\infty\cosh(\theta\rho)\,e^{-\theta^2/2}\Bigl(\frac1\theta+\frac{2c^2}{\theta^3}\Bigr)d\theta .
\]
\end{lemma}

\begin{proof}
Since $\theta^2s/2-\theta^2/2=-\theta^2c^2/2$, $\Ghat(\rho,s)=2\int_1^\infty\cosh(\theta c\rho)\,e^{-\theta^2c^2/2}(\theta^{-1}+2\theta^{-3})\,d\theta$. As $s<1$, $c>0$, and the change of variable from $\theta$ to $c\theta$ maps $[1,\infty)$ onto $[c,\infty)$, keeps $d\theta/\theta$ and turns $d\theta/\theta^3$ into $c^2d\theta/\theta^3$. This is \lname{RegretKappa.Upper.Ghat_subst} in the formalization.
\end{proof}

Splitting this integral at $\theta=1$ and subtracting $\Gamma(\rho)$ gives, since $2c^2-2=-2s$, $\Ghat(\rho,s)-\Gamma(\rho)=D(\rho,s)-4sH(\rho)$, where
\[
D(\rho,s)=2\int_c^1\cosh(\theta\rho)\,e^{-\theta^2/2}\Bigl(\frac1\theta+\frac{2c^2}{\theta^3}\Bigr)d\theta,\qquad H(\rho)=\int_1^\infty\frac{\cosh(\theta\rho)\,e^{-\theta^2/2}}{\theta^3}\,d\theta:
\]
the positive term $D$ comes from the moved cutoff, and $-4sH$ from the factor $2c^2$ in place of $2$. This gives the bound of Proposition~\ref{prop:P} for $s\le2/3$.

\begin{lemma}[small steps]\label{lem:small}
\leanmap{RegretKappa.UpperSharp.small_step_sharp, RegretKappa.UpperSharp.quartic_le}
\mapnote{Lean proves it in the normalization $\Gamma/2$ ($2.8857/2$), with the maximum of $f$ bounded by $4.32855$ from rational coefficients, and without expanding $\cosh$ in a series, from polynomial bounds on $\cosh$ with a remainder, which lead to the same quartic $f$.}
For $s\in[0,2/3]$ and $\rho\in\R$, with $\Gamma$ as in \eqref{eq:Gamma} and $\Ghat$ as in Corollary~\ref{cor:mix}, $\Ghat(\rho,s)-\Gamma(\rho)\le2.8857$.
\end{lemma}

\begin{proof}[Proof of Lemma~\ref{lem:small}]
For $s=0$, $\Ghat=\Gamma$; let $0<s\le2/3$, so that $c^2\in[1/3,1)$. By Section~\ref{subsec:small}, $\Ghat(\rho,s)-\Gamma(\rho)=D(\rho,s)-4sH(\rho)$. Expanding $\cosh$,
\begin{gather*}
\Ghat(\rho,s)-\Gamma(\rho)=\sum_{n\ge0}\frac{\rho^{2n}}{(2n)!}\bigl(D_n(s)-4sH_n\bigr),\\
D_n(s)=2\int_c^1\theta^{2n}e^{-\theta^2/2}\Bigl(\frac1\theta+\frac{2c^2}{\theta^3}\Bigr)d\theta,\qquad H_n=\int_1^\infty\theta^{2n-3}e^{-\theta^2/2}\,d\theta .
\end{gather*}
On $[c,1]$ we have $e^{-\theta^2/2}\le e^{-c^2/2}$. The function $s\mapsto-\log(1-s)$ is convex and $\tau=\frac32\log3$ is its chord slope on $[0,2/3]$, so $\log(1/c)\le\tau s/2$; and $1-c^{2m}=1-(1-s)^m\le ms$ for $m\ge1$. Hence:
\begin{itemize}
\item $D_0\le2e^{-c^2/2}\bigl(\log(1/c)+1-c^2\bigr)\le2e^{-1/6}(\tau/2+1)\,s$;
\item $D_1\le2e^{-c^2/2}\bigl((1-c^2)/2+2c^2\log(1/c)\bigr)\le2s\,e^{-c^2/2}(1/2+\tau c^2)$, and the function $v\mapsto e^{-v/2}(1/2+\tau v)$ has derivative $e^{-v/2}(\tau-1/4-\tau v/2)>0$ on $[0,1]$, so $D_1\le2e^{-1/2}(1/2+\tau)\,s$;
\item for $n\ge2$, $D_n\le2e^{-c^2/2}\bigl(\frac{1-c^{2n}}{2n}+2c^2\frac{1-c^{2n-2}}{2n-2}\bigr)\le2s\,e^{-c^2/2}(1/2+c^2)\le3e^{-1/2}s$, since $v\mapsto e^{-v/2}(1/2+v)$ has derivative $e^{-v/2}(3/4-v/2)>0$ on $[0,1]$.
\end{itemize}
Let $f(v)=\sum_{n=0}^4f_nv^n/(2n)!$ with
\[
f_0=2e^{-1/6}\Bigl(1+\frac{\tau}2\Bigr)-4H_0,\quad f_1=2e^{-1/2}\Bigl(\frac12+\tau\Bigr)-4H_1,\quad f_n=3e^{-1/2}-4H_n\ (n=2,3,4).
\]
So $(D_n(s)-4sH_n)/s\le f_n$ for $n\le4$, and for $n\ge5$ it is at most $3e^{-1/2}-4H_n\le0$, because $H_n=K_{n-1}$ increases with $n$ by Lemma~\ref{lem:Gamma}(b) and $H_2=e^{-1/2}$. Dropping the terms $n\ge5$ gives $\Ghat-\Gamma\le s\,f(\rho^2)$. By Lemma~\ref{lem:Gamma}(b), $H_1=K_0$, $H_2=e^{-1/2}$, $H_3=3e^{-1/2}$, $H_4=13e^{-1/2}$, and an integration by parts gives $H_0=(e^{-1/2}-K_0)/2$; so $f_2=-e^{-1/2}$, $f_3=-9e^{-1/2}$, $f_4=-49e^{-1/2}$, and $f(v)\le4.32855$ for every $v\ge0$ (Computation~\ref{comp:upper}). Hence $\Ghat-\Gamma\le s\,f(\rho^2)\le\frac23\cdot4.32855=2.8857$, as $0<s\le2/3$. This is \lname{RegretKappa.UpperSharp.small_step_sharp} in the formalization.
\end{proof}

\subsection{Large steps}\label{subsec:large}
For large steps, Corollary~\ref{cor:mix}(b) leaves a bound on the increase of $\Gamma$ along one move.

\begin{lemma}[large steps]\label{lem:large}
\leanmap{RegretKappa.UpperSharp.large_step_sharp, RegretKappa.UpperSharp.psi_le, RegretKappa.UpperSharp.piece_bound}
\mapnote{Lean proves it in the normalization $\Gamma/2$ ($2.7053/2$), with one theorem for each piece (\lname{RegretKappa.UpperSharp.piece_1} to \lname{RegretKappa.UpperSharp.piece_22}).}
For $s\in[2/3,1]$ and $\rho\in\R$, with $\Gamma$ as in \eqref{eq:Gamma}, $c=\sqrt{1-s}$ and $\sigma=\sqrt s$, $\max\bigl(\Gamma(\rho c+\sigma),\Gamma(\rho c-\sigma)\bigr)\le\Gamma(\rho)+2.7053$.
\end{lemma}

\begin{proof}[Proof of Lemma~\ref{lem:large}]
If $\rho\ge0$, then $|\rho c-\sigma|\le\rho c+\sigma$; if $\rho<0$, then $\Gamma(\rho c\pm\sigma)=\Gamma(|\rho|c\mp\sigma)$. Since $\Gamma$ is even and nondecreasing in $|\rho|$ (Lemma~\ref{lem:Gamma}(a)), the maximum is $\Gamma(|\rho|c+\sigma)$ in both cases, and it suffices to show $\Gamma(\rho c+\sigma)\le\Gamma(\rho)+2.7053$ for $\rho\ge0$. Let $\rho\ge0$ and $\rho_*=1/\sqrt2$.

(i) The function $s\mapsto\rho\sqrt{1-s}+\sqrt s$ is concave on $[0,1]$, with maximum $\sqrt{1+\rho^2}$ at $s=1/(1+\rho^2)$. So on $[2/3,1]$ it is at most $m(\rho)=\sqrt{1+\rho^2}$ if $\rho\le\rho_*$ (then $1/(1+\rho^2)\ge2/3$), and at most its value at $s=2/3$, $m(\rho)=\rho/\sqrt3+\sqrt{2/3}$, if $\rho\ge\rho_*$. Both expressions equal $\sqrt{3/2}$ at $\rho_*$, and $m$ is nondecreasing. Since $\Gamma$ is nondecreasing in $|\rho|$, it suffices to bound $\Gamma(m(\rho))-\Gamma(\rho)$ by $2.7053$.

(ii) Let $q(\rho)=m(\rho)^2-\rho^2$. Then $q=1$ on $[0,\rho_*]$, and $q(\rho)=-\frac23\rho^2+\frac{2\sqrt2}3\rho+\frac23$ on $[\rho_*,\infty)$, a concave quadratic with vertex at $\rho_*$; so $q$ is nonincreasing on $[0,\infty)$. If $q(\rho)\le0$, then $m(\rho)\le\rho$ and $\Gamma(m(\rho))-\Gamma(\rho)\le0$; this holds for $\rho\ge2$, since $q$ vanishes at $(1+\sqrt3)/\sqrt2<2$.

(iii) If $q(\rho)>0$, Lemma~\ref{lem:Gamma}(d) with $v=\rho^2$ and $v'=m(\rho)^2$ gives $\Gamma(m(\rho))-\Gamma(\rho)\le q(\rho)\,\delta(m(\rho)^2)$. If $\rho$ lies in a piece $[\rho^{(i-1)},\rho^{(i)}]$, this is at most $q_i\,\delta(m_i)$ for all numbers $q_i\ge q(\rho^{(i-1)})$ and $m_i\ge m(\rho^{(i)})^2$, since $q$ is nonincreasing, $m$ and $\delta$ are nondecreasing, and $\delta\ge0$ (Lemma~\ref{lem:Gamma}(c)).

(iv) Take $\rho^{(i)}=i/11$ for $i=0,\dots,22$, so that the pieces cover $[0,2]$. For $0<v<32$, Lemma~\ref{lem:Gamma}(c) and the geometric series give
\[
\delta(v)\le\sum_{n=1}^{30}\frac{n\Gamma_nv^{n-1}}{(2n)!}+\frac{2v^{30}}{31!\,(1-v/32)},
\]
with $\Gamma_n$ from Lemma~\ref{lem:Gamma}(b). With $q_i$ and $m_i$ the upward roundings to six decimals of $q(\rho^{(i-1)})$ and $m(\rho^{(i)})^2$, as in the formal proof, all $m_i$ below $32$, the twenty-two products of $q_i$ and this bound at $m_i$, with the $\Gamma_n$ computed from the rational upper bounds of $K_0$ and $e^{-1/2}$ of Computation~\ref{comp:upper}, which can only increase them, are at most $2.7053$. This is \lname{RegretKappa.UpperSharp.large_step_sharp} in the formalization.
\end{proof}

\subsection{Proofs of Proposition~\ref{prop:P} and Lemma~\ref{lem:terminal}}\label{subsec:proofP}

\begin{proof}[Proof of Proposition~\ref{prop:P}]
If $s\le2/3$, then $s<1$, and Corollary~\ref{cor:mix}(a) and Lemma~\ref{lem:small} bound the left side of \eqref{eq:P} by $2\log(\lambda+\Gamma(\rho)+2.8857)$. If $s\ge2/3$, Corollary~\ref{cor:mix}(b) and Lemma~\ref{lem:large} bound it by $2\log(\lambda+\Gamma(\rho)+2.7053)$. In both cases it is at most $2\log(\lambda+3+\Gamma(\rho))$. The split at $s=2/3$ keeps both bounds below $3$: the bound of Lemma~\ref{lem:small} is the split point times the maximum of a quartic, at most $4.32855$, and grows with the split point, while the range of steps of Lemma~\ref{lem:large}, and with it its bound, shrinks as the split point grows. This is \lname{RegretKappa.UpperSharp.step_paper} in the formalization.
\end{proof}

\begin{proof}[Proof of Lemma~\ref{lem:terminal}]
Let $\phi(z)=e^{-z^2/2}/\sqrt{2\pi}$, let $F$ be the standard normal distribution function and $\bar F=1-F$. For $\rho\ge2$,
\[
e^{-\rho^2/2}\Gamma(\rho)=\int_1^\infty e^{-(\theta-\rho)^2/2}\bigl(1+e^{-2\theta\rho}\bigr)\Bigl(\frac1\theta+\frac2{\theta^3}\Bigr)d\theta\ge\int_1^\infty\frac{e^{-(\theta-\rho)^2/2}}\theta\,d\theta .
\]
Let $M_0=\int_1^\infty e^{-(\theta-\rho)^2/2}d\theta=\sqrt{2\pi}\,F(\rho-1)$ and $M_1=\int_1^\infty\theta e^{-(\theta-\rho)^2/2}d\theta=\rho M_0+e^{-(\rho-1)^2/2}$. For $\theta>0$, $\theta^{-1}\ge2M_0/M_1-(M_0/M_1)^2\theta$, since $\theta$ times the difference is $(1-\theta M_0/M_1)^2$; multiplying by $e^{-(\theta-\rho)^2/2}$ and integrating over $\theta\ge1$ bounds the last integral below by $2M_0^2/M_1-M_0^2/M_1=M_0^2/M_1$. Hence
\[
\frac{\sqrt{2\pi}}{e^{-\rho^2/2}\Gamma(\rho)}\le\frac{\sqrt{2\pi}\,M_1}{M_0^2}=\frac{\rho}{F(\rho-1)}+\frac{\phi(\rho-1)}{F(\rho-1)^2}.
\]
For $\rho\ge2$, $\rho-1\ge1$, and $\rho/F(\rho-1)-\rho=\rho\bar F(\rho-1)/F(\rho-1)\le\frac{\rho}{\rho-1}\frac{\phi(\rho-1)}{F(\rho-1)}\le\frac{2\phi(1)}{F(1)}$, by the Mills ratio bound $\bar F(z)\le\phi(z)/z$ for $z>0$, by $\rho/(\rho-1)\le2$, and because $\phi$ decreases and $F$ increases on $[1,\infty)$; for the same reason $\phi(\rho-1)/F(\rho-1)^2\le\phi(1)/F(1)^2$. Since $2\phi(1)/F(1)+\phi(1)/F(1)^2<1$ (Computation~\ref{comp:upper}), we get $\sqrt{2\pi}/(e^{-\rho^2/2}\Gamma(\rho))\le\rho+1$, which is the inequality. For the consequence: if $|\rho|\le2$, then $C\lambda_0=e^2\ge e^{\rho^2/2}$; if $2\le|\rho|\le\sqrt T$, then $C\,\Gamma(\rho)\ge(|\rho|+1)\Gamma(\rho)/\sqrt{2\pi}\ge e^{\rho^2/2}$ by evenness. This is \lname{RegretKappa.UpperSharp.terminal_sharp} in the formalization.
\end{proof}

The numerical inequalities of this appendix are collected in the following computation.

\begin{computation}[the constants of the upper bound]\label{comp:upper}
\leanmap{RegretKappa.UpperSharp.exp_neg_half_bounds,
  RegretKappa.UpperSharp.K0_ge,
  RegretKappa.UpperSharp.K0_le,
  RegretKappa.UpperSharp.quartic_le,
  RegretKappa.UpperSharp.piece_1,
  RegretKappa.UpperSharp.mills_num,
  RegretKappa.UpperSharp.const_c_le}
\mapnote{one theorem for each piece, \lname{RegretKappa.UpperSharp.piece_1} to \lname{RegretKappa.UpperSharp.piece_22}; the rational bounds of the constants are in the module Consts}
The proofs of this appendix use the following inequalities between explicit numbers, built from $e^{-1/2}$, $e^{-1/6}$, $e^2$, $\log3$, $\pi$, $\sqrt2$, $K_0=\int_1^\infty\theta^{-1}e^{-\theta^2/2}\,d\theta$ and the standard normal density $\phi$ and distribution function $F$ at $1$:
\begin{enumerate}
\item[(1)] $0.2798867\le K_0\le0.2798869$, which gives the $\Gamma_n$ of Lemma~\ref{lem:Gamma}(b) from $e^{-1/2}$, with $0.6065306597\le e^{-1/2}\le0.6065306598$;
\item[(2)] $f(v)\le4.32855$ for $v\ge0$, for the quartic $f$ of the proof of Lemma~\ref{lem:small}, so that $\frac23\max f\le2.8857$;
\item[(3)] the twenty-two piece bounds of the proof of Lemma~\ref{lem:large}, all at most $2.7053$;
\item[(4)] $2\phi(1)/F(1)+\phi(1)/F(1)^2<1$ (Lemma~\ref{lem:terminal});
\item[(5)] $4e^{-1/2}/3\le0.81$ and $2\sqrt{2\pi}\,e^2/3\le12.35$ (Theorem~\ref{thm:upper}).
\end{enumerate}
Lean proves each of them from rational bounds on these constants proved in Lean: (1) gives two of these bounds, (2) is an inequality for a polynomial in one variable, and (3) to (5) are inequalities between explicit numbers. Item (2) is \lname{RegretKappa.UpperSharp.quartic_le} in the formalization.
\end{computation}

\section{Proofs of Sections~\ref{sec:lower} and~\ref{sec:proofmain}}\label{app:lower}

This appendix proves Lemmas~\ref{lem:chain} to~\ref{lem:phase2} and part (b) of Theorem~\ref{thm:lower}, with the notation of Section~\ref{sec:lower}: $j_0$, $L$, $k_0$, $c_1$ and $b(T)$ as in \eqref{eq:bT}, $T_1=\lfloor T/2\rfloor$, $a=\sqrt L$, $a_+=a+\sqrt{5/8}\,e^{-L/4}$, the move rule $s$ and $U(\rho)=8e^{\rho^2/2}$.

\subsection{Phase 1: a hitting time}\label{subsec:phase1}
This subsection proves Lemmas~\ref{lem:chain} and~\ref{lem:phase1}.

\begin{proof}[Proof of Lemma~\ref{lem:chain}]
By the definition of $x_{t+1}$, $V_{t+1}=V_t+V_ts_L(\rho_t)/(1-s_L(\rho_t))=V_t/(1-s_L(\rho_t))$, so $x_{t+1}^2/V_{t+1}=s_L(\rho_t)$, and $x_{t+1}\ge0$. Lemma~\ref{lem:state}(b) gives the update of $\rho$, and $y_{t+1}=\xi_{t+1}$ is a fair sign independent of the past. Also $V_1=1$ and $V_t$ does not decrease. This is \lname{RegretKappa.LowerSharp.rhoS_eq_chainOf} in the formalization.
\end{proof}

Lemma~\ref{lem:phase1} rests on a drift function. The move rule is made so that $U(\rho)=8e^{\rho^2/2}$ gains at least $1$ in expectation at each move; this drift drives $\rho^2$ to the level $L$ (Lemma~\ref{lem:FL}).

\begin{lemma}[drift]\label{lem:drift}
\leanmap{RegretKappa.LowerSharp.drift, RegretKappa.LowerSharp.driftGap_large, RegretKappa.LowerSharp.chi_le, RegretKappa.LowerSharp.drift_table}
For every $\rho\in\R$, with $s=s(\rho)$ the move rule of Section~\ref{subsec:adversary},
\[
\tfrac12\Bigl[U\bigl(\rho\sqrt{1-s}+\sqrt s\bigr)+U\bigl(\rho\sqrt{1-s}-\sqrt s\bigr)\Bigr]\ge U(\rho)+1 .
\]
\end{lemma}

Part (b) of the proof below rests on the following finite computation.

\begin{computation}[the drift certificate]\label{comp:drift}
\leanmap{RegretKappa.LowerSharp.drift_table, RegretKappa.LowerSharp.check_sound}
\mapnote{one kernel check for each cell, \lname{RegretKappa.LowerSharp.drift_cell0} to \lname{RegretKappa.LowerSharp.drift_cell27}, on the subdivisions written by \texttt{drift\_cells.gp}}
\codemap{code/drift-cells/drift_cells.gp > code/drift-cells/out/drift_cells.txt}
For $N\ge1$ let
\[
e_N(v)=\sum_{n\le N}\frac{v^n}{n!},\qquad\bar e_N(v)=e_N(v)+\frac{v^{N+1}(N+2)}{(N+1)!\,(N+2-v)},\qquad\ch_N(v)=\sum_{n\le N}\frac{v^n}{(2n)!},
\]
and for a cell index $j\in\{0,\dots,27\}$, with $s_j$ the value of the move rule on $\{|\rho|\in[j/10,(j+1)/10)\}$, and rationals $\ell<\ell'$ in $[j/10,(j+1)/10]$, let
\begin{equation}\label{eq:driftnum}
8\,e_N\Bigl(\frac{\ell^2(1-s_j)+s_j}2\Bigr)\ch_N\bigl(\ell^2s_j(1-s_j)\bigr)-8\,\bar e_N\Bigl(\frac{\ell'^2}2\Bigr)-1,
\end{equation}
a rational number that bounds from below the gain of Lemma~\ref{lem:drift} minus $1$ on $[\ell,\ell']$ (proof of Lemma~\ref{lem:drift} below). The finite statement is: for each $j=0,\dots,27$ there are rationals $j/10=\ell_0<\ell_1<\dots<\ell_n=(j+1)/10$ such that \eqref{eq:driftnum} with $N=14$, $\ell=\ell_{i-1}$ and $\ell'=\ell_i$ is nonnegative for every $i$. The subdivisions come from bisecting each cell until this holds. Lean's kernel evaluates these rational numbers in exact arithmetic, and Lean derives the inequality of Lemma~\ref{lem:drift} on $[0,2.8)$ from them. This is \lname{RegretKappa.LowerSharp.drift_table} in the formalization.
\end{computation}

\begin{proof}[Proof of Lemma~\ref{lem:drift}]
Since $(\rho\sqrt{1-s}\pm\sqrt s)^2=\rho^2(1-s)+s\pm2\rho\sqrt{s(1-s)}$, the left side is $8h_s(\rho)$ with
\[
h_s(\rho)=e^{(\rho^2(1-s)+s)/2}\cosh\bigl(\rho\sqrt{s(1-s)}\bigr).
\]
Both sides are even in $\rho$, so let $\rho\ge0$.

(a) Let $\rho\ge2.8$, so $s=\frac58e^{-\rho^2/2}$. Write $h_s(\rho)=e^{\rho^2/2}e^{-\zeta_1}\cosh\sqrt{2\zeta_2}$ with $\zeta_1=s(\rho^2-1)/2\ge0$ and $\zeta_2=\rho^2s(1-s)/2$. From $\cosh\sqrt{2\zeta_2}\ge1+\zeta_2$ and $e^{-\zeta_1}\ge1-\zeta_1\ge0$, as $\zeta_1\le\frac5{16}\rho^2e^{-\rho^2/2}<1$, we get $e^{-\zeta_1}\cosh\sqrt{2\zeta_2}\ge(1-\zeta_1)(1+\zeta_2)$, and $(1-\zeta_1)(1+\zeta_2)-1=\zeta_2-\zeta_1-\zeta_1\zeta_2$ with $\zeta_2-\zeta_1=s/2-\rho^2s^2/2$ and $\zeta_1\zeta_2=s^2\rho^2(\rho^2-1)(1-s)/4\le s^2\rho^4/4$. Hence
\[
8h_s(\rho)-U(\rho)\ge8e^{\rho^2/2}s\Bigl(\frac12-s\Bigl(\frac{\rho^2}2+\frac{\rho^4}4\Bigr)\Bigr)=\frac52-\frac{25}8\chi(\rho),\qquad\chi(\rho)=e^{-\rho^2/2}\Bigl(\frac{\rho^2}2+\frac{\rho^4}4\Bigr).
\]
With $u=\rho^2/2\ge3.92$, $\chi(\rho)=e^{-u}(u+u^2)$, and $e^u\ge\sum_{i\le5}u^i/i!\ge\frac{25}{12}(u+u^2)$, the last inequality because $120-130u-190u^2+20u^3+5u^4+u^5\ge0$ for $u\ge3.92$; so $\chi(\rho)\le12/25$. So $8h_s(\rho)-U(\rho)\ge\frac52-\frac{25}8\cdot\frac{12}{25}=1$.

(b) Let $0\le\rho<2.8$. On a cell $[j/10,(j+1)/10)$ the step $s=s_j$ is a constant rational number in $(0,1)$, and $h_s$ is nondecreasing on $[0,\infty)$, both of its factors being so. Hence on a subinterval $[\ell,\ell']$ of the cell, $8h_s(\rho)-U(\rho)\ge8h_s(\ell)-8e^{\ell'^2/2}$. With $e_N$, $\bar e_N$ and $\ch_N$ as in Computation~\ref{comp:drift}, $e_N(v)\le e^v$ for $v\ge0$, $\ch_N(v)\le\cosh\sqrt v$ for $v\ge0$, and $e^v\le\bar e_N(v)$ for $0\le v<N+2$, since the tail of the exponential series is at most its first term times a geometric series of ratio $v/(N+2)$. So $8h_s(\rho)-U(\rho)-1$ is at least the rational number \eqref{eq:driftnum} on $[\ell,\ell']$, and Computation~\ref{comp:drift} gives, with $N=14$, a subdivision of each of the $28$ cells into subintervals on which this number is nonnegative. This is \lname{RegretKappa.LowerSharp.drift} in the formalization.
\end{proof}

The drift alone does not bound the hitting time: the chain must also stay where $U$ is bounded.

\begin{lemma}[the overshoot]\label{lem:overshoot}
\leanmap{RegretKappa.LowerSharp.abs_sst_le, RegretKappa.LowerSharp.abs_rhoS_le,
  RegretKappa.LowerSharp.exp_aP_sq_le, RegretKappa.LowerSharp.U_le_G}
Let $L\ge14.5$, $a=\sqrt L$, $a_+=a+\sqrt{5/8}\,e^{-L/4}$ and $U(\rho)=8e^{\rho^2/2}$. If $|\rho|\le a_+$, both moves of the stopped chain of Lemma~\ref{lem:chain} from $\rho$ lead to points of $[-a_+,a_+]$, and $U(\rho)\le G:=8.8\,e^{L/2}$.
\end{lemma}

\begin{proof}
If $\rho^2\ge L$ the chain stays at $\rho$. Otherwise $|\rho|<a$, and a move leads to $\rho'$ with $|\rho'|\le|\rho|\sqrt{1-s}+\sqrt s\le|\rho|+\sqrt{s(\rho)}$. If $|\rho|<2.8$ this is less than $2.8+0.8<\sqrt{14.5}\le a$, as $s(\rho)\le s_0=0.6242$. If $2.8\le|\rho|<a$, it is $|\rho|+\sqrt{5/8}\,e^{-\rho^2/4}$, and $v\mapsto v+\sqrt{5/8}\,e^{-v^2/4}$ is increasing, because its derivative is at least $1-\sqrt{5/8}\max_v\frac v2e^{-v^2/4}=1-\sqrt{5/8}\,e^{-1/2}/\sqrt2>0$; so it is at most $a+\sqrt{5/8}\,e^{-a^2/4}=a_+$. Since $U$ increases with $|\rho|$, $U(\rho)\le8e^{a_+^2/2}$ for $|\rho|\le a_+$, and $a_+^2/2=L/2+\sqrt{5L/8}\,e^{-L/4}+\frac5{16}e^{-L/2}$, where the sum of the last two terms is at most $\frac45\cdot\frac19+\frac5{16}\cdot\frac1{36^2}\le0.09\le\log1.1$: indeed $\sqrt{5/8}\le4/5$, and $e^{L/4}\ge e^{29/8}\bigl(1+(L-14.5)/4\bigr)\ge36+9(L-14.5)$, which is at least $36$ and at least $9\sqrt L$, as $\sqrt L\le4+(L-14.5)$ for $L\ge14.5$. The first claim is \lname{RegretKappa.LowerSharp.abs_sst_le} in the formalization.
\end{proof}

Together, the drift and the bound on $U$ give a geometric tail for the time the chain takes to reach the level $L$, which is part (1) of Lemma~\ref{lem:phase1}.

\begin{lemma}[a geometric tail]\label{lem:FL}
\leanmap{RegretKappa.LowerSharp.lyap_add_mul_leS, RegretKappa.LowerSharp.ex_U_le, RegretKappa.LowerSharp.ex_below_pt, RegretKappa.LowerSharp.ex_below_pow}
Let $L\ge14.5$, $G=8.8\,e^{L/2}$ and $a_+=\sqrt L+\sqrt{5/8}\,e^{-L/4}$, and let $(\rho_i)_{i\ge0}$ be the stopped chain of Lemma~\ref{lem:chain} started at some $\rho_0$ with $|\rho_0|\le a_+$. For all integers $n\ge1$ and $j\ge0$, $\Prob(\rho_{jn}^2<L)\le(G/n)^j$.
\end{lemma}

\begin{proof}
By Lemma~\ref{lem:overshoot} the chain stays in $[-a_+,a_+]$, where $0\le U\le G$. While $\rho_l^2<L$ the chain moves with the step $s(\rho_l)$, so Lemma~\ref{lem:drift} gives $\E[U(\rho_{l+1})\mid\rho_l]\ge U(\rho_l)+1$; once $\rho_l^2\ge L$ it stays. Hence $\E\,U(\rho_n)\ge U(\rho_0)+\E\,\#\{l<n:\rho_l^2<L\}$. Since the chain stays once $\rho^2\ge L$, the event $\{\rho_n^2<L\}$ implies $\rho_l^2<L$ for all $l\le n$, so $\#\{l<n:\rho_l^2<L\}\ge n\mathbf 1\{\rho_n^2<L\}$, and $n\Prob(\rho_n^2<L)\le\E\,U(\rho_n)\le G$. This holds for every starting point in $[-a_+,a_+]$, and the probability is $0$ when $\rho_0^2\ge L$; so $\Prob_\rho(\rho_n^2<L)\le(G/n)\mathbf 1\{\rho^2<L\}$ for $|\rho|\le a_+$. By the Markov property at time $jn$, $\Prob(\rho_{(j+1)n}^2<L)=\E\,\Prob_{\rho_{jn}}(\rho_n^2<L)\le(G/n)\Prob(\rho_{jn}^2<L)$, and induction on $j$ concludes. This is \lname{RegretKappa.LowerSharp.ex_below_pow} in the formalization.
\end{proof}

\begin{proof}[Proof of Lemma~\ref{lem:phase1}]
By Lemma~\ref{lem:chain}, $\rho_1,\dots,\rho_{T_1}$ are $T_1-1$ steps of the stopped chain started at $\rho_1=\pm1$, and $|\rho_1|=1\le a_+$. (1) Let $G=8.8e^{L/2}$ and $n=\lceil eG\rceil\le eG+1$; then $j_0n\le T_1-1$ by (C2). Since the chain stays once $\rho^2\ge L$, $\Prob(\rho_{T_1}^2<L)$ is at most the probability that $\rho^2<L$ after $j_0n$ steps, which is at most $(G/n)^{j_0}\le e^{-j_0}$ by Lemma~\ref{lem:FL}. (2) This is Lemma~\ref{lem:overshoot} and the definition of $\Succ$. (3) Fix $t\le T_1$. The prediction $\yh_t$ is a function of $\xi_1,\dots,\xi_{t-1}$, because the features and outcomes before round $t$ and the feature $x_t$ are. The outcome $y_t=\xi_t$ is independent of them, so $\E(\yh_t^2-2\yh_ty_t)=\E\,\yh_t^2\ge0$. Sum over $t\le T_1$. Part (1) is \lname{RegretKappa.LowerSharp.prob_fail_le} in the formalization.
\end{proof}

\subsection{Phase 2: the van Trees inequality with the past as a prior}\label{subsec:phase2}
This subsection proves Lemmas~\ref{lem:split} and~\ref{lem:phase2}.

\begin{proof}[Proof of Lemma~\ref{lem:split}]
Let $V=V_{T_1}$ and $\bar y=\sum_i\alpha_iy'_i$. By Lemma~\ref{lem:chain}, $S_{T_1}=\rho\sqrt V$, so $S_T=\sqrt V(\rho+\bar y)$ and $V_T=V(1+Q)$, and $\rho_T^2=(\rho+\bar y)^2/(1+Q)$. Expanding the squares gives $\sum_i[(\yh'_i-y'_i)^2-((\rho+z)\alpha_i-y'_i)^2]=\sum_i(\yh_i'^2-2\yh'_iy'_i)-(\rho+z)^2Q+2(\rho+z)\bar y$ and $(1+Q)((\rho+\bar y)/(1+Q)-\rho-z)^2=\rho_T^2-2(\rho+z)(\rho+\bar y)+(\rho+z)^2(1+Q)$, so $\rho^2+\Delta(z)=\sum_i(\yh_i'^2-2\yh'_iy'_i)+\rho_T^2+\rho^2-z^2-2(\rho+z)\rho+(\rho+z)^2$, and the last four terms cancel. Lemma~\ref{lem:identity} concludes. Off $\Succ$, the rounds after $T_1$ have $x_t=y_t=0$, so $\rho_T=\rho_{T_1}$ and, by Lemma~\ref{lem:identity}, $\Reg(T)=\sum_{t\le T_1}(\yh_t^2-2\yh_ty_t)+\sum_{t>T_1}\yh_t^2+\rho_{T_1}^2$. The identity on $\Succ$ is \lname{RegretKappa.LowerSharp.regret_succ} in the formalization.
\end{proof}

The proof of Lemma~\ref{lem:phase2} rests on the van Trees inequality \citep[Section~2, inequality~(3)]{GL}, in the following form for finitely many signs.

\begin{lemma}[van Trees inequality]\label{lem:vT}
\leanmap{RegretKappa.LowerSharp.vanTreesG}
\mapnote{Lean proves it for the probabilities $\frac12(1+y_l(a_l+b_lz))$ with $|a_l+b_lz|\le c_l<1$ and the bound $1/(J_g+\sum_lb_l^2/(1-c_l^2))$, of which the lemma is the case $a_l=\rho\alpha_l$, $b_l=\alpha_l$, $c_l=\nu_l$.}
Let $g\colon\R\to\R$ be continuously differentiable with $g(-2)=g(2)=0$, $\int_{-2}^2g^2=1$ and $J_g=\int_{-2}^24g'^2>0$. Let $n\ge0$, and let $\rho$, $\alpha_l$ and $\nu_l$ be real numbers with $\nu_l<1$ and $|(\rho+z)\alpha_l|\le\nu_l$ for $z\in[-2,2]$ and $l=1,\dots,n$. For $y\in\{-1,1\}^n$ and $z\in[-2,2]$ let $P_z(y)=\prod_l\frac12\bigl(1+y_l(\rho+z)\alpha_l\bigr)$. Then for every function $\psi\colon\{-1,1\}^n\to\R$,
\[
\int_{-2}^2\sum_{y}g(z)^2\,P_z(y)\,\bigl(\psi(y)-z\bigr)^2dz\ge\frac1{J_g+\sum_l\alpha_l^2/(1-\nu_l^2)}.
\]
\end{lemma}

\begin{proof}
The weights $P_z(y)$ are positive, since $\nu_l<1$. For each $y$, integration by parts gives $\int_{-2}^2(\psi(y)-z)\,\partial_z\bigl(g^2P_z(y)\bigr)\,dz=\int_{-2}^2g^2P_z(y)\,dz$, the boundary terms vanishing because $g(\pm2)=0$; summing over $y$, since $\sum_yP_z(y)=1$ and $\int g^2=1$, $\sum_y\int(\psi(y)-z)\,\partial_z(g^2P_z(y))\,dz=1$. Now $\partial_z(g^2P_z(y))=g\,P_z(y)\bigl(g\,\varsigma+2g'\bigr)$ with the score $\varsigma(z,y)=\sum_ly_l\alpha_l/(1+y_l(\rho+z)\alpha_l)$. By the Cauchy-Schwarz inequality for the measure $P_z(y)\,dz$ times counting measure,
\[
1\le\Bigl[\sum_y\int g^2P_z(y)(\psi(y)-z)^2dz\Bigr]\Bigl[\sum_y\int P_z(y)\bigl(g\,\varsigma+2g'\bigr)^2dz\Bigr].
\]
Under $P_z$ the signs are independent, and $y_l\alpha_l/(1+y_l(\rho+z)\alpha_l)$ has mean $0$ and second moment $\alpha_l^2/(1-((\rho+z)\alpha_l)^2)\le\alpha_l^2/(1-\nu_l^2)$. So $\sum_yP_z\varsigma=0$ and $\sum_yP_z\varsigma^2\le\sum_l\alpha_l^2/(1-\nu_l^2)$, and the second factor is at most $\int\bigl(g^2\sum_l\alpha_l^2/(1-\nu_l^2)+4g'^2\bigr)dz=\sum_l\alpha_l^2/(1-\nu_l^2)+J_g$. This is \lname{RegretKappa.LowerSharp.vanTreesG} in the formalization.
\end{proof}

The proof of Lemma~\ref{lem:phase2} uses the following constants of the prior of the adversary.

\begin{lemma}[the prior]\label{lem:prior}
\leanmap{RegretKappa.LowerSharp.integral_gP_sq, RegretKappa.LowerSharp.integral_gP_deriv_sq, RegretKappa.LowerSharp.integral_sq_gP_sq, RegretKappa.LowerSharp.gP_two, RegretKappa.LowerSharp.gP_neg_two}
The function $g(z)=\cos(\pi z/4)/\sqrt2$ of Section~\ref{subsec:adversary}, whose square is the density of $Z$, is continuously differentiable on $\R$, with $g(\pm2)=0$, $\int_{-2}^2g^2=1$, $\int_{-2}^24g'^2=J=\pi^2/4$ and $\int_{-2}^2z^2g(z)^2\,dz=\frac43-\frac8{\pi^2}$.
\end{lemma}

\begin{proof}
With $z=2v$, $\int_{-2}^2g^2=\int_{-1}^1\cos^2(\pi v/2)\,dv=1$. Next $g'(z)=-\frac\pi{4\sqrt2}\sin(\pi z/4)$, so $\int_{-2}^24g'^2=\frac{\pi^2}8\int_{-2}^2\sin^2(\pi z/4)\,dz=\frac{\pi^2}4$. For the second moment, $\int_{-2}^2z^2g^2=\frac14\int_{-2}^2z^2(1+\cos(\pi z/2))\,dz=\frac14\bigl(\frac{16}3-\frac{32}{\pi^2}\bigr)$, since $\int_{-1}^1v^2\cos(\pi v)\,dv=-\frac2\pi\int_{-1}^1v\sin(\pi v)\,dv=-\frac4{\pi^2}$. The value of $J$ is \lname{RegretKappa.LowerSharp.integral_gP_deriv_sq} in the formalization.
\end{proof}

\begin{proof}[Proof of Lemma~\ref{lem:phase2}]
Condition on the phase-1 signs throughout. Then $\rho$, $r$ and all phase-2 features are fixed numbers, $Z$ has the density $g^2$ of Lemma~\ref{lem:prior}, and given $Z=z$ the outcomes $y'_1,\dots,y'_{k_0}$ are independent signs with $\Prob(y'_l=1)=\frac12(1+(\rho+z)\alpha_l)$: the model of Lemma~\ref{lem:vT}, since $|(\rho+z)\alpha_l|\le r\alpha_l=\nu_l<1$.

(i) Given $Z$ and $y'_1,\dots,y'_{i-1}$, the prediction $\yh'_i$ is fixed and $y'_i$ has mean $(\rho+Z)\alpha_i$, so the conditional mean of $(\yh'_i-y'_i)^2-((\rho+Z)\alpha_i-y'_i)^2$ is $(\yh'_i-(\rho+Z)\alpha_i)^2$. Also $\E Z^2=\frac43-\frac8{\pi^2}$ by Lemma~\ref{lem:prior}.

(ii) Round $T_1+i$. The prediction $\yh'_i$ is a function of $y'_1,\dots,y'_{i-1}$; write $\yh'_i=\alpha_i(\psi+\rho)$, which defines $\psi$, as $\alpha_i>0$. Then $(\yh'_i-(\rho+Z)\alpha_i)^2=\alpha_i^2(\psi-Z)^2$, and Lemma~\ref{lem:vT} for the first $i-1$ outcomes, with $J_g=J$, gives $\E(\psi-Z)^2\ge1/(J+\sum_{l<i}\alpha_l^2/(1-\nu_l^2))=1/\Info_{i-1}$. Hence $\E(\yh'_i-(\rho+Z)\alpha_i)^2\ge\alpha_i^2/\Info_{i-1}$.

(iii) The number $(\rho+\sum_i\alpha_iy'_i)/(1+Q)-\rho$ is a function of the $k_0$ outcomes, and Lemma~\ref{lem:vT} with all of them gives $\E[(1+Q)(\dots-\rho-Z)^2]\ge(1+Q)/\Info_{k_0}$.

Summing (i) to (iii) gives the first claim.

(iv) Let $\zeta_i=\log(\Info_i/J)$, so $\zeta_0=0$. Since $\alpha_i^2=(1-\nu_i^2)(\Info_i-\Info_{i-1})$ and $v\ge\log(1+v)$, we have $\alpha_i^2/\Info_{i-1}\ge(1-\nu_i^2)(\zeta_i-\zeta_{i-1})$. With $\nu_i^2=i/(2k_0)$, summation by parts gives
\[
\sum_{i=1}^{k_0}\frac{\alpha_i^2}{\Info_{i-1}}\ge\frac12\,\zeta_{k_0}+\frac1{2k_0}\sum_{i=0}^{k_0-1}\zeta_i .
\]
Put $\Lambda=\log(k_0/(\pi^2r^2))$, which is at most $\log k_0$. Since $\Info_i\ge J+\sum_{l\le i}\alpha_l^2=J\bigl(1+e^\Lambda i(i+1)/k_0^2\bigr)$, we have $\zeta_i\ge\log(1+e^\Lambda i^2/k_0^2)\ge0$. Then $\frac12\zeta_{k_0}\ge\Lambda/2$, and, $\log(1+e^\Lambda v^2/k_0^2)$ being increasing in $v\ge0$,
\[
\frac1{2k_0}\sum_{i=1}^{k_0-1}\log\Bigl(1+\frac{e^\Lambda i^2}{k_0^2}\Bigr)\ge\frac1{2k_0}\int_0^{k_0-1}\log\frac{e^\Lambda v^2}{k_0^2}\,dv=\frac{k_0-1}{2k_0}\Bigl(\Lambda+2\log\Bigl(1-\frac1{k_0}\Bigr)-2\Bigr).
\]
Now $2\log(1-1/k_0)\ge-2/(k_0-1)$, so the right side is at least $\frac{k_0-1}{2k_0}(\Lambda-2)-\frac1{k_0}\ge\frac{\Lambda-2}2-\frac\Lambda{2k_0}-\frac1{k_0}$. Hence $\sum_i\alpha_i^2/\Info_{i-1}\ge\Lambda-1-\log(k_0)/(2k_0)-1/k_0$. For the last term, $\Info_{k_0}\le J+2Q$ because $\nu_i^2\le1/2$, and $Q=(k_0+1)/(4r^2)$, so
\[
\frac{1+Q}{\Info_{k_0}}\ge\frac{1+Q}{J+2Q}=\frac12-\frac{J/2-1}{J+2Q}\ge\frac12-\frac{J-2}{4Q}\ge\frac12-\frac{J-2}{\pi^2e^2},
\]
the last step because $4Q=(k_0+1)/r^2\ge\pi^2e^2$ when $k_0\ge\pi^2e^2r^2$; this is the only use of that condition. Adding, $\beta(\rho)\ge\log(k_0/r^2)-\log(\pi^2)-1-\E Z^2+\frac12-\frac{J-2}{\pi^2e^2}-\frac{\log k_0}{2k_0}-\frac1{k_0}=\underline\beta(\rho)$. The first claim is \lname{RegretKappa.LowerSharp.phase2S} in the formalization.
\end{proof}

\subsection{The side conditions and part (b) of Theorem~\ref{thm:lower}}
This subsection proves Lemma~\ref{lem:side} and part (b) of Theorem~\ref{thm:lower}.

\begin{proof}[Proof of Lemma~\ref{lem:side}]
(C1) If $j_0\le4$, then $L\ge2\log(T_0/(80e))\ge14.5$, since $T_0\ge80\cdot4446\ge80\,e^{33/4}$. If $j_0=j\ge5$, then $\log(3\log T)>j-1$ by the definition of $j_0$, so $2\log T>\frac23e^{j-1}$, and $L=2\log T-2\log(20ej)>\frac23e^{j-1}-2\log(20ej)\ge14.5$: indeed $2\log(20ej)\le2j+6$, as $\log20\le3$ and $\log j\le j-1$, and $\frac23e^{j-1}\ge\frac23e^4(j-4)\ge36(j-4)\ge2j+20.5$ for $j\ge5$, since $e^{j-5}\ge j-4$ and $e^4\ge54.5$. (C2) By the definition of $L$, $8.8\,e\,e^{L/2}j_0=0.44\,T$, and $0.44\,T+j_0\le T/2-3/2\le T_1-1$ because $0.06\,T\ge j_0+3/2$ for $T\ge T_0$, as $j_0<\log(3\log T)+1\le\log\log T+2.1$. (C3) Since $L\le2\log T$ and $\sqrt{5/8}\,e^{-L/4}\le1$, $(a_++2)^2\le(\sqrt{2\log T}+3)^2\le4\log T+18$, and $\pi^2e^2\le74$. As $\log T\le\log T_0+(T-T_0)/T_0\le13+(T-T_0)/T_0$, $74(4\log T+18)\le5180+296(T-T_0)/T_0\le T_0/2+(T-T_0)/2=T/2\le k_0$. Condition (C1) is \lname{RegretKappa.LowerSharp.level_ge} in the formalization.
\end{proof}

\begin{proof}[Proof of Theorem~\ref{thm:lower}, part (b)]
Here $L=2\log T-2\log(20e)-2\log j_0$ with $j_0\le\log\log T+2.1$, as in the proof of (C2); $\log k_0\ge\log(T/2)$; $2\log(a+2)=\log L+2\log(1+2/a)\le\log\log T+\log2+0.85$, since $L\le2\log T$ and $a\ge3.8$, so that $2\log(1+2/a)\le2\log(29/19)\le0.85$; $\log(k_0)/(2k_0)+1/k_0\le10^{-4}$; and $c_1\le3.33$. Collecting, and since $2\log(20e)+2\log2+0.85+3.33+10^{-4}\le13.59$,
\[
b_0\ge3\log T-\log\log T-2\log(\log\log T+2.1)-13.6 ,
\]
with $b_0$ as in the proof of part (a). Next let $z=\log\log T$. For $T\ge T_0$, $\log T\ge12.19\ge e^{2.5}$, so $z\ge2.5$, and with $\log v\le v-1$ and $2\log4.6\le3.1$,
\[
2\log(z+2.1)\le2\log4.6+\frac{2(z-2.5)}{4.6}\le z+0.6 .
\]
So the right side is at least $3\log T-2\log\log T-14.2$, which is nonnegative, as $\log\log T\le\log T-1$ and $\log T\ge12.2$; in particular $b_0\ge0$. Since $e^{-j_0}\le1/(3\log T)$ and $b_0\le3\log T$, $b(T)=(1-e^{-j_0})b_0\ge b_0-1$, which gives the first form, and the bound $2\log(z+2.1)\le z+0.6$ gives the second. This is \lname{RegretKappa.LowerSharp.lowerSharpSimplified} in the formalization.
\end{proof}

\subsection{Proof of Corollary~\ref{cor:bounded}}

\begin{proof}[Proof of Corollary~\ref{cor:bounded}]
Let $(x^i,y^i,p_i)$ be the adversary of Theorem~\ref{thm:lower}(a), and $X=1+\sum_{i,t}x^i_t$, which depends on $T$ only. The new adversary plays $(x^i/X,By^i)$ with probability $p_i$; its features lie in $[0,1]$, since $x^i_t\ge0$, and its outcomes in $\{-B,0,B\}$. For a deterministic learner $\mathcal L$, let $\mathcal L'$ predict $\frac1B\mathcal L_t(x_1/X,\dots,x_t/X;By_1,\dots,By_{t-1})$; this is a learner, because $X$ depends on $T$ only. On $(x^i,y^i)$ its predictions are $1/B$ times those of $\mathcal L$ on $(x^i/X,By^i)$, and $\inf_\theta\sum_t(\theta x^i_t/X-By^i_t)^2=B^2\inf_\theta\sum_t(\theta x^i_t-y^i_t)^2$ by the substitution $\theta\mapsto XB\theta$; so the regret of $\mathcal L$ on $(x^i/X,By^i)$ is $B^2$ times the regret of $\mathcal L'$ on $(x^i,y^i)$, and $\sum_ip_i\Reg_{\mathcal L}(x^i/X,By^i)\ge B^2b(T)$ by Theorem~\ref{thm:lower}(a). For a randomized learner, this holds for each $\mathcal L^\omega$; adding $\sum_ip_i\sum_t(By^i_t)^2$ and integrating against the probability $\mu$ gives expected regret at least $B^2b(T)$, which is at least $B^2(3\log T-2\log\log T-15.2)$ by Theorem~\ref{thm:lower}(b). In the chain, the first inequality holds because, for each deterministic learner, some play of the new adversary has regret at least the weighted average; the second because $\{-B,0,B\}\subseteq[-B,B]$; and the third by the learner of Theorem~\ref{thm:upper}, whose regret is at most $B^2u(T)\le B^2(3\log T+0.37)$. This is \lname{RegretKappa.mainBoundedFeatures} in the formalization.
\end{proof}

\bibliographystyle{plainnat}
\bibliography{refs}

\end{document}
